% Options for packages loaded elsewhere
\PassOptionsToPackage{unicode}{hyperref}
\PassOptionsToPackage{hyphens}{url}
\documentclass[
]{article}
\usepackage{xcolor}
\usepackage{amsmath,amssymb}
\setcounter{secnumdepth}{-\maxdimen} % remove section numbering
\usepackage{iftex}
\ifPDFTeX
  \usepackage[T1]{fontenc}
  \usepackage[utf8]{inputenc}
  \usepackage{textcomp} % provide euro and other symbols
\else % if luatex or xetex
  \usepackage{unicode-math} % this also loads fontspec
  \defaultfontfeatures{Scale=MatchLowercase}
  \defaultfontfeatures[\rmfamily]{Ligatures=TeX,Scale=1}
\fi
\usepackage{lmodern}
\ifPDFTeX\else
  % xetex/luatex font selection
\fi
% Use upquote if available, for straight quotes in verbatim environments
\IfFileExists{upquote.sty}{\usepackage{upquote}}{}
\IfFileExists{microtype.sty}{% use microtype if available
  \usepackage[]{microtype}
  \UseMicrotypeSet[protrusion]{basicmath} % disable protrusion for tt fonts
}{}
\makeatletter
\@ifundefined{KOMAClassName}{% if non-KOMA class
  \IfFileExists{parskip.sty}{%
    \usepackage{parskip}
  }{% else
    \setlength{\parindent}{0pt}
    \setlength{\parskip}{6pt plus 2pt minus 1pt}}
}{% if KOMA class
  \KOMAoptions{parskip=half}}
\makeatother
\usepackage{color}
\usepackage{fancyvrb}
\newcommand{\VerbBar}{|}
\newcommand{\VERB}{\Verb[commandchars=\\\{\}]}
\DefineVerbatimEnvironment{Highlighting}{Verbatim}{commandchars=\\\{\}}
% Add ',fontsize=\small' for more characters per line
\newenvironment{Shaded}{}{}
\newcommand{\AlertTok}[1]{\textcolor[rgb]{1.00,0.00,0.00}{\textbf{#1}}}
\newcommand{\AnnotationTok}[1]{\textcolor[rgb]{0.38,0.63,0.69}{\textbf{\textit{#1}}}}
\newcommand{\AttributeTok}[1]{\textcolor[rgb]{0.49,0.56,0.16}{#1}}
\newcommand{\BaseNTok}[1]{\textcolor[rgb]{0.25,0.63,0.44}{#1}}
\newcommand{\BuiltInTok}[1]{\textcolor[rgb]{0.00,0.50,0.00}{#1}}
\newcommand{\CharTok}[1]{\textcolor[rgb]{0.25,0.44,0.63}{#1}}
\newcommand{\CommentTok}[1]{\textcolor[rgb]{0.38,0.63,0.69}{\textit{#1}}}
\newcommand{\CommentVarTok}[1]{\textcolor[rgb]{0.38,0.63,0.69}{\textbf{\textit{#1}}}}
\newcommand{\ConstantTok}[1]{\textcolor[rgb]{0.53,0.00,0.00}{#1}}
\newcommand{\ControlFlowTok}[1]{\textcolor[rgb]{0.00,0.44,0.13}{\textbf{#1}}}
\newcommand{\DataTypeTok}[1]{\textcolor[rgb]{0.56,0.13,0.00}{#1}}
\newcommand{\DecValTok}[1]{\textcolor[rgb]{0.25,0.63,0.44}{#1}}
\newcommand{\DocumentationTok}[1]{\textcolor[rgb]{0.73,0.13,0.13}{\textit{#1}}}
\newcommand{\ErrorTok}[1]{\textcolor[rgb]{1.00,0.00,0.00}{\textbf{#1}}}
\newcommand{\ExtensionTok}[1]{#1}
\newcommand{\FloatTok}[1]{\textcolor[rgb]{0.25,0.63,0.44}{#1}}
\newcommand{\FunctionTok}[1]{\textcolor[rgb]{0.02,0.16,0.49}{#1}}
\newcommand{\ImportTok}[1]{\textcolor[rgb]{0.00,0.50,0.00}{\textbf{#1}}}
\newcommand{\InformationTok}[1]{\textcolor[rgb]{0.38,0.63,0.69}{\textbf{\textit{#1}}}}
\newcommand{\KeywordTok}[1]{\textcolor[rgb]{0.00,0.44,0.13}{\textbf{#1}}}
\newcommand{\NormalTok}[1]{#1}
\newcommand{\OperatorTok}[1]{\textcolor[rgb]{0.40,0.40,0.40}{#1}}
\newcommand{\OtherTok}[1]{\textcolor[rgb]{0.00,0.44,0.13}{#1}}
\newcommand{\PreprocessorTok}[1]{\textcolor[rgb]{0.74,0.48,0.00}{#1}}
\newcommand{\RegionMarkerTok}[1]{#1}
\newcommand{\SpecialCharTok}[1]{\textcolor[rgb]{0.25,0.44,0.63}{#1}}
\newcommand{\SpecialStringTok}[1]{\textcolor[rgb]{0.73,0.40,0.53}{#1}}
\newcommand{\StringTok}[1]{\textcolor[rgb]{0.25,0.44,0.63}{#1}}
\newcommand{\VariableTok}[1]{\textcolor[rgb]{0.10,0.09,0.49}{#1}}
\newcommand{\VerbatimStringTok}[1]{\textcolor[rgb]{0.25,0.44,0.63}{#1}}
\newcommand{\WarningTok}[1]{\textcolor[rgb]{0.38,0.63,0.69}{\textbf{\textit{#1}}}}
\usepackage{longtable,booktabs,array}
\usepackage{caption}
\captionsetup[table]{skip=6pt}
\newcounter{none} % for unnumbered tables
\usepackage{calc} % for calculating minipage widths
% Correct order of tables after \paragraph or \subparagraph
\usepackage{etoolbox}
\makeatletter
\patchcmd\longtable{\par}{\if@noskipsec\mbox{}\fi\par}{}{}
\makeatother
% Allow footnotes in longtable head/foot
\IfFileExists{footnotehyper.sty}{\usepackage{footnotehyper}}{\usepackage{footnote}}
\makesavenoteenv{longtable}
\setlength{\emergencystretch}{3em} % prevent overfull lines
\providecommand{\tightlist}{%
  \setlength{\itemsep}{0pt}\setlength{\parskip}{0pt}}
\usepackage{bookmark}
\IfFileExists{xurl.sty}{\usepackage{xurl}}{} % add URL line breaks if available
\urlstyle{same}
% fallback for those not using the hyperref driver hyperxmp:
\makeatletter
\@ifundefined{xmpquote}{\newcommand{\xmpquote}[1]{#1}}{}
\makeatother
\hypersetup{
  hidelinks,
  pdfcreator={LaTeX via pandoc}}

\author{}
\date{}

\begin{document}

{
\setcounter{tocdepth}{3}
\tableofcontents
}
\section{TÀI LIỆU HỌC TẬP}\label{tuxe0i-liux1ec7u-hux1ecdc-tux1eadp}

\section{CƠ SỞ HẠ TẦNG KHÓA CÔNG KHAI (PKI) \& CHUYỂN DỊCH MẬT MÃ HẬU
LƯỢNG TỬ
(PQC)}\label{cux1a1-sux1edf-hux1ea1-tux1ea7ng-khuxf3a-cuxf4ng-khai-pki-chuyux1ec3n-dux1ecbch-mux1eadt-muxe3-hux1eadu-lux1b0ux1ee3ng-tux1eed-pqc}

\textbf{Mã học phần:} ATHK 512 · \textbf{Khóa đào tạo:} Cao học An toàn
Thông tin (CHAT4P)\\
\textbf{Học kỳ:} HK261 · \textbf{Năm học:} 2026 -- 2027\\
\textbf{Cơ sở đào tạo:} Khoa An toàn Thông tin --- Học viện Kỹ thuật Mật
mã (Phân hiệu TP.HCM)\\
\textbf{Đơn vị phối hợp chuyên môn:} Khoa Khoa học \& Kỹ thuật Máy tính,
Trường Đại học Bách khoa --- ĐHQG TP.HCM

\begin{center}\rule{0.5\linewidth}{0.5pt}\end{center}

\subsubsection{THÔNG TIN ĐỘI NGŨ GIẢNG DẠY \& HỖ TRỢ HỌC
THUẬT}\label{thuxf4ng-tin-ux111ux1ed9i-ngux169-giux1ea3ng-dux1ea1y-hux1ed7-trux1ee3-hux1ecdc-thuux1eadt}

\begin{itemize}
\item
  \textbf{Chủ biên \& Giảng viên phụ trách:}\\
  \textbf{Nguyễn An Khương}\\
  \emph{Đơn vị công tác:} Khoa Khoa học \& Kỹ thuật Máy tính, Trường Đại
  học Bách khoa --- ĐHQG TP.HCM\\
  \emph{Mã định danh học thuật toàn cầu:}
  \href{https://orcid.org/0000-0002-9910-6387}{ORCID:
  0000-0002-9910-6387}\\
  \emph{Email công vụ:} \texttt{nakhuong@gmail.com} /
  \texttt{nakhuong@hcmut.edu.vn}
\item
  \textbf{Trợ giảng học phần (Teaching Assistant):}\\
  \textbf{Nguyễn Hữu Hoàng Huy}\\
  \emph{Hồ sơ chuyên môn:}
  \href{https://www.linkedin.com/in/nguyen-huu-hoang-huy-abb138366/}{LinkedIn
  Profile}\\
  \emph{GitHub:} \texttt{@HuyLHP} (Quyền quản trị kỹ thuật: Collaborator
  \texttt{write})
\end{itemize}

\begin{center}\rule{0.5\linewidth}{0.5pt}\end{center}

\subsubsection{BẢN QUYỀN VÀ GIẤY PHÉP PHÁT
HÀNH}\label{bux1ea3n-quyux1ec1n-vuxe0-giux1ea5y-phuxe9p-phuxe1t-huxe0nh}

\begin{itemize}
\tightlist
\item
  \textbf{Bản quyền Nội dung \& Học liệu:} Bản quyền thuộc về Giảng viên
  \textbf{Nguyễn An Khương}. Tác phẩm này được phát hành dưới Giấy phép
  Quốc tế
  \href{https://creativecommons.org/licenses/by-nc-nd/4.0/}{Creative
  Commons Attribution-NonCommercial-NoDerivatives 4.0 (CC BY-NC-ND
  4.0)}. Người học được phép sao chép, phân phối tài liệu cho mục đích
  học tập cá nhân phi thương mại, không được phép chỉnh sửa hoặc cắt
  ghép nội dung khi chưa có văn bản chấp thuận của tác giả.
\item
  \textbf{Bản quyền Mã nguồn \& Thực nghiệm:} Toàn bộ mã nguồn, kịch bản
  tự động hóa, cấu hình container Docker và test suites trong học liệu
  được phát hành theo \href{https://opensource.org/licenses/MIT}{Giấy
  phép Nguồn mở MIT (MIT License)}.
\end{itemize}

\begin{center}\rule{0.5\linewidth}{0.5pt}\end{center}

\subsection{LỜI NÓI ĐẦU}\label{lux1eddi-nuxf3i-ux111ux1ea7u}

Hạ tầng Khóa Công khai (Public Key Infrastructure - PKI) từ lâu đã trở
thành ``xương sống vô hình'' bảo đảm tính bảo mật, toàn vẹn, xác thực và
chống chối bỏ cho toàn bộ thế giới số hiện đại. Từ giao thức bắt tay TLS
bảo vệ hàng tỷ phiên duyệt web an toàn mỗi ngày, chữ ký số trên các
chứng từ tài chính - ngân hàng, đến các định danh thiết bị IoT và cơ chế
khởi động an toàn (Secure Boot) trong hệ điều hành máy tính --- tất cả
đều vận hành dựa trên các nguyên lý cốt lõi của PKI.

Tuy nhiên, kỷ nguyên số đang bước vào một khúc quanh lịch sử với hai
chuyển biến mang tính cách mạng: 1. \textbf{Kiến trúc Zero Trust \& Điện
toán Đám mây Phân tán:} Mô hình chu vi mạng truyền thống bị xóa bỏ hoàn
toàn. Các hệ thống hiện đại đòi hỏi hàng triệu kết nối mTLS giữa các
microservices được mã hóa với chứng chỉ có thời hạn sống biểu kiến
(ephemeral certificates tính bằng giờ hoặc phút), kéo theo yêu cầu tự
động hóa tuyệt đối thông qua giao thức ACME (RFC 8555) và kiểm định tuân
thủ tự động (ZLint). 2. \textbf{Nguy cơ Hiện hữu từ Điện toán Lượng tử
(Post-Quantum Transition):} Thuật toán lượng tử Shor đã chứng minh khả
năng bẻ gãy hoàn toàn các hệ mật mã khóa công khai kinh điển dựa trên
bài toán phân tích thừa số nguyên tố (RSA) và logarit rời rạc (DSA,
ECDSA, ECDH) trong thời gian đa thức. Tháng 8/2024, Viện Tiêu chuẩn và
Công nghệ Quốc gia Hoa Kỳ (NIST) đã chính thức ban hành bộ tiêu chuẩn
FIPS 203 (ML-KEM), FIPS 204 (ML-DSA) và FIPS 205 (SLH-DSA), mở ra một
cuộc đại di chuyển mật mã (cryptographic migration) quy mô toàn cầu chưa
từng có trong lịch sử nhân loại.

Bộ \textbf{Tài liệu học tập: Cơ sở Hạ tầng Khóa công khai \& Mật mã Hậu
lượng tử} này được biên soạn dành riêng cho học viên cao học ngành An
toàn Thông tin (CHAT4P) nhằm trang bị một góc nhìn sâu sắc, toàn diện và
hiện đại: kết hợp chặt chẽ giữa \textbf{nền tảng toán học hàn lâm chuẩn
mực}, \textbf{chuẩn giao thức kỹ thuật quốc tế (IETF RFCs, CA/B Forum,
NIST FIPS)} và \textbf{năng lực thực chiến công nghiệp (DevSecOps,
Docker, OpenSSL 3.5, liboqs, HSM)}.

Tài liệu được kết cấu thành 8 chương chuyên sâu kèm theo hệ thống phụ
lục thực nghiệm phong phú, đóng vai trò là kim chỉ nam xuyên suốt 9 buổi
học lý thuyết, thực hành phòng lab và 5 đề tài đồ án chuyên đề lớn của
học phần.

\begin{center}\rule{0.5\linewidth}{0.5pt}\end{center}

\subsection{MỤC LỤC TỔNG QUAN}\label{mux1ee5c-lux1ee5c-tux1ed5ng-quan}

\begin{itemize}
\tightlist
\item
  \textbf{Chương 1:} Nền Tảng Lý Thuyết Số \& Cơ Sở Toán Học Mật Mã Khóa
  Công Khai
\item
  \textbf{Chương 2:} Cấu Trúc Dữ Liệu Chứng Chỉ X.509 v3 \& Quy Cách Mã
  Hóa ASN.1 DER
\item
  \textbf{Chương 3:} Kiến Trúc Phân Tầng CA, Vòng Đời Khóa \& Quy Chế
  CP/CPS (RFC 3647)
\item
  \textbf{Chương 4:} Thuật Toán Xác Thực Đường Dẫn Tin Cậy \& Cơ Chế Thu
  Hồi Chứng Chỉ
\item
  \textbf{Chương 5:} Giao Thức Quản Trị Tự Động ACME (RFC 8555) \& Hạ
  Tầng Zero Trust PKI
\item
  \textbf{Chương 6:} Các Kịch Bản Tấn Công Hạ Tầng PKI \& Chiến Lược
  Phòng Thủ Toàn Diện
\item
  \textbf{Chương 7:} Chuyển Dịch Hệ Thống Sang Mật Mã Hậu Lượng Tử (PQC
  Migration)
\item
  \textbf{Chương 8:} Khung Pháp Lý, Quy Chuẩn Kỹ Thuật \& Hạ Tầng Chữ Ký
  Số Quốc Gia
\item
  \textbf{Phụ lục A:} Cẩm Nang 50 Lệnh OpenSSL 3.5 Thực Chiến Cho Chuyên
  Gia ATTT
\item
  \textbf{Phụ lục B:} Bộ Công Cụ Python Cryptography Phân Tích \& Kiểm
  Định Chứng Chỉ
\item
  \textbf{Phụ lục C:} Danh Mục Tài Liệu Tham Khảo Hàn Lâm \& Khóa Học
  Quốc Tế (OCW)
\end{itemize}

\begin{center}\rule{0.5\linewidth}{0.5pt}\end{center}

\section{CHƯƠNG 1: NỀN TẢNG LÝ THUYẾT SỐ \& CƠ SỞ TOÁN HỌC MẬT MÃ KHÓA
CÔNG
KHAI}\label{chux1b0ux1a1ng-1-nux1ec1n-tux1ea3ng-luxfd-thuyux1ebft-sux1ed1-cux1a1-sux1edf-touxe1n-hux1ecdc-mux1eadt-muxe3-khuxf3a-cuxf4ng-khai}

\subsection{1.1. Cấu Trúc Đại Số Trừu Tượng \& Các Phép Toán Số
Học}\label{cux1ea5u-truxfac-ux111ux1ea1i-sux1ed1-trux1eebu-tux1b0ux1ee3ng-cuxe1c-phuxe9p-touxe1n-sux1ed1-hux1ecdc}

Mật mã khóa công khai không thể tồn tại nếu thiếu nền tảng vững chắc của
đại số trừu tượng. Thay vì hoạt động trên tập số thực liên tục
\(\mathbb{R}\), mật mã học hiện đại vận hành trên các cấu trúc rời rạc
hữu hạn nhằm bảo đảm tính chính xác tuyệt đối và ngăn chặn triệt để hiện
tượng sai số làm tròn số học.

\subsubsection{1.1.1. Nhóm (Groups), Vành (Rings) và Trường Hữu Hạn
(Finite
Fields)}\label{nhuxf3m-groups-vuxe0nh-rings-vuxe0-trux1b0ux1eddng-hux1eefu-hux1ea1n-finite-fields}

\begin{enumerate}
\def\labelenumi{\arabic{enumi}.}
\item
  \textbf{Nhóm (Group):} Một tập hợp \(G\) cùng với phép toán hai ngôi
  \(\star\) được gọi là một nhóm \((G, \star)\) nếu thỏa mãn 4 tiên đề:

  \begin{itemize}
  \tightlist
  \item
    \textbf{Tính đóng (Closure):}
    \(\forall a, b \in G \implies a \star b \in G\).
  \item
    \textbf{Tính kết hợp (Associativity):}
    \(\forall a, b, c \in G \implies (a \star b) \star c = a \star (b \star c)\).
  \item
    \textbf{Phần tử trung hòa (Identity):}
    \(\exists e \in G, \forall a \in G \implies a \star e = e \star a = a\).
  \item
    \textbf{Phần tử nghịch đảo (Inverse):}
    \(\forall a \in G, \exists a^{-1} \in G \implies a \star a^{-1} = a^{-1} \star a = e\).
  \item
    Nếu \(a \star b = b \star a\), nhóm được gọi là \textbf{nhóm giao
    hoán (nhóm Abel)}.
  \end{itemize}
\item
  \textbf{Vành (Ring):} Một tập hợp \(R\) cùng hai phép toán cộng
  \((+)\) và nhân \((\cdot)\) được gọi là một vành nếu \((R, +)\) là một
  nhóm Abel, phép nhân có tính kết hợp và phân phối đối với phép cộng.
  Vành số dư \(\mathbb{Z}_n = \{0, 1, \dots, n-1\}\) là vành cơ bản của
  hệ mật RSA.
\item
  \textbf{Trường (Field):} Một vành giao hoán có đơn vị mà trong đó mọi
  phần tử khác 0 đều khả nghịch dưới phép nhân. Trường hữu hạn Galois
  \(\mathbb{F}_q\) chỉ tồn tại khi số phần tử \(q = p^m\) (với \(p\) là
  số nguyên tố và \(m \ge 1\)):

  \begin{itemize}
  \tightlist
  \item
    Khi \(m = 1\), ta có trường số nguyên tố
    \(\mathbb{F}_p = \mathbb{Z}_p\), nơi các phép toán modulo \(p\) diễn
    ra.
  \item
    Khi \(m > 1\), ta có trường mở rộng \(\mathbb{F}_{2^m}\), đặc biệt
    quan trọng trong phần cứng mật mã vi mạch.
  \end{itemize}
\end{enumerate}

\subsubsection{1.1.2. Các Định Lý Số Học Cốt
Lõi}\label{cuxe1c-ux111ux1ecbnh-luxfd-sux1ed1-hux1ecdc-cux1ed1t-luxf5i}

\begin{itemize}
\item
  \textbf{Bổ đề Bézout \& Thuật toán Euclid mở rộng:} Với hai số nguyên
  \(a, b\), luôn tồn tại hai số nguyên \(x, y\) sao cho:
  \[a \cdot x + b \cdot y = \gcd(a, b)\] Khi \(\gcd(a, n) = 1\), ta có
  \(a \cdot x \equiv 1 \pmod n\), tức \(x = a^{-1} \pmod n\). Thuật toán
  Euclid mở rộng tính toán phần tử nghịch đảo này trong độ phức tạp thời
  gian cực nhanh \(O(\log n)\).
\item
  \textbf{Định lý Fermat nhỏ \& Định lý Euler:}

  \begin{itemize}
  \tightlist
  \item
    Nếu \(p\) nguyên tố và \(\gcd(a, p) = 1\), thì
    \(a^{p-1} \equiv 1 \pmod p\).
  \item
    Tổng quát hóa bởi Leonhard Euler: Với số nguyên dương \(n\), hàm phi
    Euler \(\phi(n)\) đếm số lượng các số nguyên dương nhỏ hơn \(n\) và
    nguyên tố cùng nhau với \(n\). Nếu \(n = p \cdot q\) (\(p, q\)
    nguyên tố), thì: \[\phi(n) = (p - 1)(q - 1)\]
  \item
    Khi đó, với mọi \(a\) nguyên tố cùng nhau với \(n\):
    \[a^{\phi(n)} \equiv 1 \pmod n\]
  \end{itemize}
\item
  \textbf{Định lý Phần dư Trung Hoa (Chinese Remainder Theorem - CRT):}
  Cho các số nguyên dương \(m_1, m_2, \dots, m_k\) đôi một nguyên tố
  cùng nhau. Hệ phương trình đồng dư:
  \[x \equiv a_i \pmod{m_i} \quad (i = 1, \dots, k)\] luôn có nghiệm duy
  nhất modulo \(M = \prod_{i=1}^k m_i\). Trong cài đặt thực tế của RSA,
  kỹ thuật RSA-CRT phân rã phép tính lũy thừa trên module 2048-bit thành
  hai phép tính trên hai module \(p\) và \(q\) 1024-bit, giúp tăng tốc
  độ giải mã và ký số lên \textbf{4 lần}.
\end{itemize}

\begin{center}\rule{0.5\linewidth}{0.5pt}\end{center}

\subsection{1.2. Khái Niệm Hàm Một Chiều \& Cửa Sập (Trapdoor One-Way
Functions)}\label{khuxe1i-niux1ec7m-huxe0m-mux1ed9t-chiux1ec1u-cux1eeda-sux1eadp-trapdoor-one-way-functions}

Mọi hệ thống mật mã khóa công khai đều được xây dựng trên nền tảng của
\textbf{Hàm một chiều có cửa sập (Trapdoor One-Way Function - TOWF)}:

\[f: X \to Y\]

Hàm \(f\) phải thỏa mãn ba đặc trưng toán học nghiêm ngặt: 1. \textbf{Dễ
tính toán theo chiều xuôi:} Cho trước \(x \in X\), việc tính
\(y = f(x)\) là một bài toán thời gian đa thức \(O(\text{poly}(|x|))\).
2. \textbf{Bất khả thi theo chiều ngược (Hard to Invert):} Cho trước
\(y \in Y\), việc tìm \(x\) sao cho \(f(x) = y\) là bài toán bất khả thi
về mặt độ phức tạp tính toán đối với bất kỳ thuật toán cổ điển nào (thời
gian tính toán vượt xa tuổi thọ vũ trụ). 3. \textbf{Cửa sập bí mật
(Trapdoor Key):} Tồn tại một thông tin bổ trợ đặc biệt
\(K_{\text{priv}}\) (khóa bí mật). Khi biết \(K_{\text{priv}}\), việc
tính ngược \(x = f^{-1}(y, K_{\text{priv}})\) trở thành một bài toán đa
thức tường minh và hiệu quả.

\begin{verbatim}
                  +-----------------------------------+
                  |          Đầu vào x                |
                  +-----------------------------------+
                                    |
                                    | Chiều xuôi: Dễ (Đa thức)
                                    v
                  +-----------------------------------+
                  |          Đầu ra y = f(x)          |
                  +-----------------------------------+
                                    |
       Không có Cửa sập             |             Có Cửa sập (K_priv)
  ---------------------------       |       -------------------------------
  Nghịch đảo: BẤT KHẢ THI           |       Nghịch đảo: DỄ DÀNG (Đa thức)
  (Chi phí hàm mũ hoặc cận hàm mũ)  v       x = f^-1(y, K_priv)
\end{verbatim}

Hiện nay, toán học thế giới công nhận ba bài toán khó kinh điển được
dùng làm nền tảng cho các hàm cửa sập trong PKI cổ điển: 1. \textbf{Bài
toán Phân tích Thừa số Nguyên tố (Integer Factorization Problem - IFP)}
\(\implies\) Hệ mật RSA. 2. \textbf{Bài toán Logarit Rời rạc (Discrete
Logarithm Problem - DLP)} \(\implies\) Giao thức Diffie-Hellman,
ElGamal, DSA. 3. \textbf{Bài toán Logarit Rời rạc trên Đường cong
Elliptic (ECDLP)} \(\implies\) ECDH, ECDSA, Ed25519.

\begin{center}\rule{0.5\linewidth}{0.5pt}\end{center}

\subsection{1.3. Bài Toán Phân Tích Thừa Số Nguyên Tố \& Hệ Mật Mã
RSA}\label{buxe0i-touxe1n-phuxe2n-tuxedch-thux1eeba-sux1ed1-nguyuxean-tux1ed1-hux1ec7-mux1eadt-muxe3-rsa}

\subsubsection{1.3.1. Thuật Toán Sinh Khóa, Mã Hóa và Ký Số
RSA}\label{thuux1eadt-touxe1n-sinh-khuxf3a-muxe3-huxf3a-vuxe0-kuxfd-sux1ed1-rsa}

Được phát minh bởi Ron Rivest, Adi Shamir và Leonard Adleman năm 1977
tại MIT: 1. \textbf{Sinh khóa:} - Chọn ngẫu nhiên hai số nguyên tố rất
lớn \(p\) và \(q\) (mỗi số 1024-bit hoặc 2048-bit). - Tính module công
khai \(n = p \cdot q\). Kích thước \(n\) thường là 2048-bit hoặc
4096-bit. - Tính hàm Euler: \(\phi(n) = (p - 1)(q - 1)\). - Chọn số mũ
công khai \(e\) thỏa mãn \(1 < e < \phi(n)\) và
\(\gcd(e, \phi(n)) = 1\). Giá trị chuẩn công nghiệp luôn chọn
\(e = 65537 = 2^{16} + 1\) (số nguyên tố Fermat \(F_4\)) nhằm tối ưu hóa
phép lũy thừa nhị phân (chỉ có đúng hai bit 1). - Tính số mũ bí mật
\(d\) bằng thuật toán Euclid mở rộng:
\(d \equiv e^{-1} \pmod{\phi(n)}\), tức
\(e \cdot d \equiv 1 \pmod{\phi(n)}\). - \textbf{Khóa công khai:}
\(K_{\text{pub}} = (n, e)\). - \textbf{Khóa bí mật:}
\(K_{\text{priv}} = (d, p, q)\).

\begin{enumerate}
\def\labelenumi{\arabic{enumi}.}
\setcounter{enumi}{1}
\tightlist
\item
  \textbf{Mã hóa \& Giải mã:}

  \begin{itemize}
  \tightlist
  \item
    Mã hóa bản rõ \(m < n\): \(c = m^e \pmod n\).
  \item
    Giải mã bản mã \(c\): \(m = c^d \pmod n\).
  \item
    Tính đúng đắn xuất phát từ Định lý Euler:
    \(c^d \equiv (m^e)^d \equiv m^{e \cdot d} \equiv m^{1 + k \cdot \phi(n)} \equiv m \cdot (m^{\phi(n)})^k \equiv m \cdot 1^k \equiv m \pmod n\).
  \end{itemize}
\item
  \textbf{Chữ ký số RSA:}

  \begin{itemize}
  \tightlist
  \item
    Người ký tính hàm băm thông điệp \(h = H(M)\).
  \item
    Chữ ký số: \(\sigma = h^d \pmod n\).
  \item
    Người kiểm tra tính: \(h' = \sigma^e \pmod n\). Nếu \(h' == H(M)\),
    chữ ký hợp lệ.
  \end{itemize}
\end{enumerate}

\subsubsection{1.3.2. Cơ Chế Đệm An Toàn (Padding Schemes) \& Các Nguy
Cơ Tấn
Công}\label{cux1a1-chux1ebf-ux111ux1ec7m-an-touxe0n-padding-schemes-cuxe1c-nguy-cux1a1-tux1ea5n-cuxf4ng}

RSA ``thô'' (Textbook RSA) hoàn toàn không an toàn do tính chất
\textbf{đồng cấu nhân}:
\[(m_1)^e \cdot (m_2)^e \equiv (m_1 \cdot m_2)^e \pmod n\] Kẻ tấn công
có thể giả mạo chữ ký trên thông điệp mới từ tích của hai chữ ký đã
biết. Do đó, chuẩn mã hóa PKI bắt buộc áp dụng cơ chế đệm ngẫu nhiên: *
\textbf{PKCS\#1 v1.5:} Thêm chuỗi padding ngẫu nhiên có định dạng cố
định. Tuy nhiên, PKCS\#1 v1.5 dễ bị tổn thương trước tấn công
Bleichenbacher (Million Message Attack) trên các máy chủ SSL/TLS nếu trả
về lỗi giải mã không đồng nhất. * \textbf{RSA-OAEP (Optimal Asymmetric
Encryption Padding):} Dùng cho mã hóa, bảo đảm tính an toàn IND-CCA2
(kháng tấn công bản mã chọn lọc thích nghi). * \textbf{RSA-PSS
(Probabilistic Signature Scheme):} Được chuẩn hóa trong RFC 8017 / FIPS
186-5. Sử dụng muối ngẫu nhiên (salt) trước khi băm, chứng minh an toàn
tuyệt đối trong mô hình Random Oracle Model. Các hệ thống PKI hiện đại
được khuyến nghị chuyển dịch toàn bộ từ RSA PKCS\#1 v1.5 sang RSA-PSS.

\begin{center}\rule{0.5\linewidth}{0.5pt}\end{center}

\subsection{1.4. Bài Toán Logarit Rời Rạc \& Mật Mã Đường Cong Elliptic
(ECC)}\label{buxe0i-touxe1n-logarit-rux1eddi-rux1ea1c-mux1eadt-muxe3-ux111ux1b0ux1eddng-cong-elliptic-ecc}

\subsubsection{1.4.1. Bài Toán Logarit Rời Rạc
(DLP)}\label{buxe0i-touxe1n-logarit-rux1eddi-rux1ea1c-dlp}

Cho nhóm cyclic \(G\) cấp \(q\) sinh bởi phần tử \(g\). - \textbf{Bài
toán chiều xuôi:} Cho \(x \in \mathbb{Z}_q\), tính \(y = g^x \in G\) (dễ
dàng bằng thuật toán nhân đôi lũy thừa). - \textbf{Bài toán Logarit Rời
rạc (DLP):} Cho \(y \in G\), tìm \(x \in \mathbb{Z}_q\) sao cho
\(g^x = y\). - Trên nhóm nhân \(\mathbb{Z}_p^*\), các thuật toán Sàng
trường số tổng quát (General Number Field Sieve - GNFS) và Index
Calculus giải quyết DLP trong thời gian cận hàm mũ:
\[L_p\left[\frac{1}{3}, c\right] = \exp\left((c + o(1)) (\ln p)^{1/3} (\ln \ln p)^{2/3}\right)\]
Do đó, để đạt mức an toàn 128-bit, module số nguyên tố \(p\) phải có độ
dài ít nhất \textbf{3072 bits}.

\subsubsection{1.4.2. Mật Mã Đường Cong Elliptic (Elliptic Curve
Cryptography -
ECC)}\label{mux1eadt-muxe3-ux111ux1b0ux1eddng-cong-elliptic-elliptic-curve-cryptography---ecc}

Nhận định thiên tài của Neal Koblitz và Victor Miller năm 1985: Nếu thay
thế nhóm nhân \(\mathbb{Z}_p^*\) bằng nhóm các điểm trên đường cong
Elliptic, các thuật toán Index Calculus hoàn toàn bị vô hiệu hóa!

\begin{enumerate}
\def\labelenumi{\arabic{enumi}.}
\item
  \textbf{Phương trình Đường cong Elliptic:} Trên trường hữu hạn
  \(\mathbb{F}_p\) (\(p > 3\)), đường cong Weierstrass không suy biến có
  phương trình: \[E: y^2 \equiv x^3 + a \cdot x + b \pmod p\] với điều
  kiện biệt thức \(\Delta = 4a^3 + 27b^2 \not\equiv 0 \pmod p\).
\item
  \textbf{Luật Cộng Điểm Hình Học \& Đại Số:} Tập hợp các điểm
  \(P = (x, y) \in E(\mathbb{F}_p)\) cùng với \textbf{Điểm ở vô cực
  \(\mathcal{O}\)} (phần tử trung hòa) tạo thành một nhóm Abel:

  \begin{itemize}
  \tightlist
  \item
    Với hai điểm \(P = (x_1, y_1)\) và \(Q = (x_2, y_2)\): Nếu
    \(P \neq Q\), hệ số góc đường cát tuyến:
    \(\lambda = \frac{y_2 - y_1}{x_2 - x_1} \pmod p\). Nếu \(P = Q\)
    (nhân đôi điểm), hệ số góc tiếp tuyến:
    \(\lambda = \frac{3x_1^2 + a}{2y_1} \pmod p\).
  \item
    Tọa độ điểm tổng \(R = P + Q = (x_3, y_3)\):
    \[x_3 = \lambda^2 - x_1 - x_2 \pmod p\]
    \[y_3 = \lambda(x_1 - x_3) - y_1 \pmod p\]
  \end{itemize}
\item
  \textbf{Ưu thế Vượt Trội Của ECC Trong Hạ Tầng PKI:} Thuật toán mạnh
  nhất hiện nay giải bài toán ECDLP trên các đường cong chuẩn (NIST
  P-256, Curve25519) chỉ là Pollard's Rho với độ phức tạp hàm mũ đầy đủ:
  \[\mathcal{O}(\sqrt{\pi \cdot n / 2}) \approx 2^{128} \text{ phép tính đối với khóa 256-bit}\]
\end{enumerate}

{\def\LTcaptype{none} % do not increment counter
\begin{longtable}[]{@{}
  >{\centering\arraybackslash}p{(\linewidth - 8\tabcolsep) * \real{0.2000}}
  >{\centering\arraybackslash}p{(\linewidth - 8\tabcolsep) * \real{0.2000}}
  >{\centering\arraybackslash}p{(\linewidth - 8\tabcolsep) * \real{0.2000}}
  >{\centering\arraybackslash}p{(\linewidth - 8\tabcolsep) * \real{0.2000}}
  >{\centering\arraybackslash}p{(\linewidth - 8\tabcolsep) * \real{0.2000}}@{}}
\toprule\noalign{}
\begin{minipage}[b]{\linewidth}\centering
Mức Độ An Toàn (Security Bits)
\end{minipage} & \begin{minipage}[b]{\linewidth}\centering
Kích Thước Khóa Đối Xứng (AES)
\end{minipage} & \begin{minipage}[b]{\linewidth}\centering
Kích Thước Khóa RSA (Bits)
\end{minipage} & \begin{minipage}[b]{\linewidth}\centering
Kích Thước Khóa ECC (Bits)
\end{minipage} & \begin{minipage}[b]{\linewidth}\centering
Tỉ Lệ Tiết Kiệm Kích Thước
\end{minipage} \\
\midrule\noalign{}
\endhead
\bottomrule\noalign{}
\endlastfoot
80-bit (Lỗi thời) & 2TDEA & 1024 & 160 & 6.4 : 1 \\
112-bit (Tối thiểu) & 3TDEA & 2048 & 224 & 9.1 : 1 \\
\textbf{128-bit (Chuẩn hiện đại)} & \textbf{AES-128} & \textbf{3072} &
\textbf{256} & \textbf{12.0 : 1} \\
\textbf{192-bit (An ninh cao)} & \textbf{AES-192} & \textbf{7680} &
\textbf{384} & \textbf{20.0 : 1} \\
\textbf{256-bit (Tuyệt mật quốc gia)} & \textbf{AES-256} &
\textbf{15360} & \textbf{521} & \textbf{29.5 : 1} \\
\end{longtable}
}

Chính sự chênh lệch vượt trội này khiến toàn bộ hạ tầng Web PKI, thiết
bị di động thông minh và các thẻ căn cước gắn chip hiện đại đều ưu tiên
chuyển dịch sang \textbf{ECDSA} và \textbf{Ed25519} (Edwards curve RFC
8032) nhằm tiết kiệm băng thông mạng và năng lượng tính toán.

\begin{center}\rule{0.5\linewidth}{0.5pt}\end{center}

\section{CHƯƠNG 2: CẤU TRÚC DỮ LIỆU CHỨNG CHỈ X.509 V3 \& QUY CÁCH MÃ
HÓA ASN.1
DER}\label{chux1b0ux1a1ng-2-cux1ea5u-truxfac-dux1eef-liux1ec7u-chux1ee9ng-chux1ec9-x.509-v3-quy-cuxe1ch-muxe3-huxf3a-asn.1-der}

\subsection{2.1. Cú Pháp ASN.1 \& Quy Tắc Mã Hóa Nhị Phân
DER}\label{cuxfa-phuxe1p-asn.1-quy-tux1eafc-muxe3-huxf3a-nhux1ecb-phuxe2n-der}

Chứng chỉ số không phải là một chuỗi văn bản ngẫu nhiên mà là một cấu
trúc dữ liệu phân cấp được định nghĩa bằng chuẩn quốc tế \textbf{ASN.1
(Abstract Syntax Notation One - ITU-T X.680)} và được đóng gói nhị phân
theo quy tắc \textbf{DER (Distinguished Encoding Rules - ITU-T X.690)}.

\subsubsection{2.1.1. Cấu Trúc Khung Nhị Phân TLV (Tag - Length -
Value)}\label{cux1ea5u-truxfac-khung-nhux1ecb-phuxe2n-tlv-tag---length---value}

Trong mã hóa DER, mỗi trường dữ liệu bắt buộc tuân theo định dạng 3
phần: 1. \textbf{Tag (Loại dữ liệu):} 1 byte (hoặc nhiều bytes) xác định
kiểu: - \texttt{0x02}: \texttt{INTEGER} (Số nguyên) - \texttt{0x03}:
\texttt{BIT\ STRING} (Chuỗi bit) - \texttt{0x04}: \texttt{OCTET\ STRING}
(Chuỗi byte) - \texttt{0x05}: \texttt{NULL} - \texttt{0x06}:
\texttt{OBJECT\ IDENTIFIER\ (OID)} - \texttt{0x0C}: \texttt{UTF8String}
- \texttt{0x13}: \texttt{PrintableString} - \texttt{0x17}:
\texttt{UTCTime} (YYMMDDHHMMSSZ) - \texttt{0x18}:
\texttt{GeneralizedTime} (YYYYMMDDHHMMSSZ) - \texttt{0x30}:
\texttt{SEQUENCE} (Cấu trúc danh sách tuần tự - bit Constructed được
bật: \texttt{0x20\ \textbar{}\ 0x10\ =\ 0x30}) - \texttt{0x31}:
\texttt{SET} (Tập hợp không thứ tự)

\begin{enumerate}
\def\labelenumi{\arabic{enumi}.}
\setcounter{enumi}{1}
\tightlist
\item
  \textbf{Length (Độ dài phần dữ liệu Value):}

  \begin{itemize}
  \tightlist
  \item
    \textbf{Dạng ngắn (Short form):} Nếu chiều dài \(\le 127\) bytes,
    trường Length chiếm đúng 1 byte với bit cao nhất bằng 0
    (\texttt{0b0xxxxxxx}). Ví dụ: dài 45 bytes \(\implies\)
    \texttt{0x2D}.
  \item
    \textbf{Dạng dài (Long form):} Nếu chiều dài \(> 127\) bytes, byte
    đầu tiên có bit cao nhất bằng 1 (\texttt{0b1xxxxxxx}), trong đó 7
    bits sau chỉ số lượng byte tiếp theo dùng để biểu diễn chiều dài. Ví
    dụ: chiều dài 500 bytes (\texttt{0x01F4}), biểu diễn:
    \texttt{0x82\ 0x01\ 0xF4}.
  \item
    \textbf{Quy tắc bất biến của DER:} Khác với BER (Basic Encoding
    Rules) cho phép độ dài không xác định (indefinite length), DER bắt
    buộc phải dùng độ dài xác định ngắn nhất có thể. Điều này bảo đảm
    tính đơn nhất (canonical representation): \textbf{Một đối tượng
    ASN.1 chỉ có duy nhất một biểu diễn nhị phân DER tương ứng}, ngăn
    chặn hoàn toàn tấn công giả mạo chữ ký bằng cách biến đổi padding
    nhị phân (malleability attack).
  \end{itemize}
\end{enumerate}

\begin{center}\rule{0.5\linewidth}{0.5pt}\end{center}

\subsection{2.2. Khung Dữ Liệu TBSCertificate Theo Chuẩn RFC
5280}\label{khung-dux1eef-liux1ec7u-tbscertificate-theo-chuux1ea9n-rfc-5280}

Cấu trúc chứng chỉ X.509 v3 được mô tả bằng cú pháp ASN.1 chính thức
trong RFC 5280:

\begin{Shaded}
\begin{Highlighting}[]
\NormalTok{Certificate ::= }\DataTypeTok{SEQUENCE}\NormalTok{ \{}
\NormalTok{    tbsCertificate          TBSCertificate,}
\NormalTok{    signatureAlgorithm      AlgorithmIdentifier,}
\NormalTok{    signatureValue          BIT }\DataTypeTok{STRING}
\NormalTok{\}}

\NormalTok{TBSCertificate ::= }\DataTypeTok{SEQUENCE}\NormalTok{ \{}
\NormalTok{    version             [}\DecValTok{0}\NormalTok{] EXPLICIT Version }\KeywordTok{DEFAULT}\NormalTok{ v1,}
\NormalTok{    serialNumber            CertificateSerialNumber,}
\NormalTok{    signature               AlgorithmIdentifier,}
\NormalTok{    issuer                  Name,}
\NormalTok{    validity                Validity,}
\NormalTok{    subject                 Name,}
\NormalTok{    subjectPublicKeyInfo    SubjectPublicKeyInfo,}
\NormalTok{    issuerUniqueID      [}\DecValTok{1}\NormalTok{] IMPLICIT UniqueIdentifier }\KeywordTok{OPTIONAL}\NormalTok{,}
\NormalTok{    subjectUniqueID     [}\DecValTok{2}\NormalTok{] IMPLICIT UniqueIdentifier }\KeywordTok{OPTIONAL}\NormalTok{,}
\NormalTok{    extensions          [}\DecValTok{3}\NormalTok{] EXPLICIT Extensions }\KeywordTok{OPTIONAL}
\NormalTok{\}}
\end{Highlighting}
\end{Shaded}

\begin{verbatim}
+-------------------------------------------------------------------+
|                     X.509 v3 CERTIFICATE                          |
|                                                                   |
| +---------------------------------------------------------------+ |
| |                    TBSCertificate (To-Be-Signed)              | |
| |                                                               | |
| | 1. Version: v3 (0x02)                                         | |
| | 2. Serial Number: Số nguyên ngẫu nhiên dương (tối đa 20 bytes)| |
| | 3. Signature Algorithm: e.g. ecdsa-with-SHA256 / sha256WithRSA| |
| | 4. Issuer DN: CN=KMA Intermediate CA, O=KMA, C=VN             | |
| | 5. Validity: notBefore (02/10/2026) -> notAfter (02/10/2027)  | |
| | 6. Subject DN: CN=api.kma.edu.vn, O=KMA, C=VN                 | |
| | 7. SubjectPublicKeyInfo: Algorithm Identifier + Public Key Bit| |
| | 8. Extensions (X.509 v3 Extensions):                          | |
| |    - Basic Constraints: CA:FALSE                              | |
| |    - Key Usage: Digital Signature                             | |
| |    - Extended Key Usage: Server Auth, Client Auth             | |
| |    - SAN: DNS:api.kma.edu.vn, IP:10.0.0.1                     | |
| |    - CRL Distribution Points: http://crl.kma.edu.vn/ca.crl    | |
| |    - AIA: OCSP http://ocsp.kma.edu.vn                         | |
| +---------------------------------------------------------------+ |
|                                                                   |
| Signature Algorithm Identifier: ecdsa-with-SHA256 (OID 1.2.840...) |
| Signature Value: Chữ ký số của CA trên toàn bộ khối DER băm của   |
|                  TBSCertificate [BIT STRING]                      |
+-------------------------------------------------------------------+
\end{verbatim}

\begin{center}\rule{0.5\linewidth}{0.5pt}\end{center}

\subsection{2.3. Phân Tích Chuyên Sâu Các Phần Mở Rộng Trọng Yếu
(Critical
Extensions)}\label{phuxe2n-tuxedch-chuyuxean-suxe2u-cuxe1c-phux1ea7n-mux1edf-rux1ed9ng-trux1ecdng-yux1ebfu-critical-extensions}

Phần mở rộng (Extensions) là nơi hiện thực hóa toàn bộ chính sách an
ninh của PKI. Mỗi phần mở rộng gồm: \texttt{extnID} (OID), cờ
\texttt{critical} (BOOLEAN, mặc định FALSE) và \texttt{extnValue} (OCTET
STRING chứa dữ liệu DER đóng gói). \textgreater{} \textbf{Quy Tắc Vàng
Của RFC 5280 Về Cờ Critical:} Nếu một phần mở rộng có cờ
\texttt{critical\ =\ TRUE}, mọi ứng dụng hoặc thiết bị xác thực (trình
duyệt, hệ điều hành) \textbf{BẮT BUỘC} phải hiểu và thực thi đúng ngữ
nghĩa của phần mở rộng đó. Nếu không nhận biết hoặc không thể xử lý, hệ
thống \textbf{BẮT BUỘC PHẢI TỪ CHỐI} chứng chỉ ngay lập tức!

\begin{enumerate}
\def\labelenumi{\arabic{enumi}.}
\tightlist
\item
  \textbf{Basic Constraints (OID \texttt{2.5.29.19} - BẮT BUỘC CRITICAL
  trên CA):}

  \begin{itemize}
  \tightlist
  \item
    Chứa trường \texttt{cA} (BOOLEAN) và \texttt{pathLenConstraint}
    (INTEGER tùy chọn).
  \item
    Nếu \texttt{cA\ =\ TRUE}: Chứng chỉ đại diện cho một nhà cầm quyền
    xác thực (CA), được phép ký cấp phát chứng chỉ khác.
  \item
    Nếu \texttt{cA\ =\ FALSE}: Chứng chỉ thực thể cuối (End-Entity /
    Leaf Certificate), tuyệt đối không được dùng để ký chứng chỉ khác.
  \item
    \texttt{pathLenConstraint\ =\ n}: Quy định chuỗi CA con bên dưới nó
    chỉ được phép có tối đa \(n\) cấp Intermediate CA trung gian trước
    khi tới chứng chỉ lá. Nếu \(n = 0\), CA này chỉ được phép ký trực
    tiếp cho End-Entity, không được ký cấp phát thêm bất kỳ Sub-CA nào
    nữa.
  \item
    \emph{Lỗ hổng thực tế:} Nếu quản trị viên sơ suất cấp chứng chỉ cho
    máy chủ nội bộ mà đặt \texttt{cA\ =\ TRUE}, kẻ tấn công chiếm được
    máy chủ đó có thể tự biến mình thành CA trung gian và phát hành
    chứng chỉ giả mạo cho mọi tên miền trên Internet!
  \end{itemize}
\item
  \textbf{Key Usage (OID \texttt{2.5.29.15} - CRITICAL):} Định nghĩa các
  thao tác mật mã mà khóa công khai trong chứng chỉ được phép thực hiện:

  \begin{itemize}
  \tightlist
  \item
    \texttt{digitalSignature}: Ký thông điệp hoặc dữ liệu (trừ chứng chỉ
    và CRL).
  \item
    \texttt{nonRepudiation}: Chống chối bỏ (thường dùng trong chữ ký văn
    bản số, pháp lý).
  \item
    \texttt{keyEncipherment}: Mã hóa khóa đối xứng (dùng trong RSA Key
    Transport).
  \item
    \texttt{keyAgreement}: Trao đổi khóa (dùng trong ECDH).
  \item
    \texttt{keyCertSign}: Ký cấp phát chứng chỉ số (Chỉ được phép bật
    nếu \texttt{cA\ =\ TRUE}).
  \item
    \texttt{cRLSign}: Ký danh sách thu hồi chứng chỉ (CRL).
  \end{itemize}
\item
  \textbf{Extended Key Usage - EKU (OID \texttt{2.5.29.37}):} Xác định
  mục đích ứng dụng cụ thể:

  \begin{itemize}
  \tightlist
  \item
    \texttt{serverAuth} (\texttt{1.3.6.1.5.5.7.3.1}): Xác thực máy chủ
    Web TLS.
  \item
    \texttt{clientAuth} (\texttt{1.3.6.1.5.5.7.3.2}): Xác thực người
    dùng / máy trạm mTLS.
  \item
    \texttt{codeSigning} (\texttt{1.3.6.1.5.5.7.3.3}): Ký mã nguồn phần
    mềm, driver.
  \item
    \texttt{timeStamping} (\texttt{1.3.6.1.5.5.7.3.8}): Cấp dấu thời
    gian RFC 3161.
  \end{itemize}
\item
  \textbf{Subject Alternative Name - SAN (OID \texttt{2.5.29.17}):}

  \begin{itemize}
  \tightlist
  \item
    Danh sách định danh thực thể thay thế: \texttt{dNSName},
    \texttt{iPAddress}, \texttt{uniformResourceIdentifier},
    \texttt{rfc822Name} (email).
  \item
    \textbf{Quy định nghiêm ngặt từ CA/B Forum Baseline Requirements
    v1.8+:} Các trình duyệt hiện đại (Chrome, Firefox, Safari) đã loại
    bỏ hoàn toàn việc đọc tên miền từ trường \texttt{Common\ Name\ (CN)}
    của Subject DN. Toàn bộ tên miền và địa chỉ IP hợp lệ \textbf{BẮT
    BUỘC PHẢI KHAI BÁO TRONG PHẦN MỞ RỘNG SAN}. Chứng chỉ chỉ có CN mà
    thiếu SAN sẽ bị trình duyệt báo lỗi
    \texttt{NET::ERR\_CERT\_COMMON\_NAME\_INVALID}.
  \end{itemize}
\end{enumerate}

\begin{center}\rule{0.5\linewidth}{0.5pt}\end{center}

\section{CHƯƠNG 3: KIẾN TRÚC PHÂN TẦNG CA, VÒNG ĐỜI KHÓA \& QUY CHẾ
CP/CPS (RFC
3647)}\label{chux1b0ux1a1ng-3-kiux1ebfn-truxfac-phuxe2n-tux1ea7ng-ca-vuxf2ng-ux111ux1eddi-khuxf3a-quy-chux1ebf-cpcps-rfc-3647}

\subsection{3.1. Các Mô Hình Phân Tầng Hạ Tầng Chứng Thực (CA
Hierarchies)}\label{cuxe1c-muxf4-huxecnh-phuxe2n-tux1ea7ng-hux1ea1-tux1ea7ng-chux1ee9ng-thux1ef1c-ca-hierarchies}

Trong thực tế, một hệ thống PKI không bao giờ triển khai dưới dạng một
CA đơn lẻ bởi vì nguyên tắc bảo vệ tài sản thông tin tối thượng:
\textbf{Càng ít sử dụng khóa bí mật có giá trị cao thì hệ thống càng an
toàn.}

\begin{verbatim}
                           +------------------------+
                           |     ROOT CA (GỐC)      |
                           |  - Tự ký (Self-signed) |
                           |  - Offline / Air-gap   |
                           |  - Khóa: 4096-bit RSA  |
                           |    hoặc P-384 / ML-DSA |
                           |  - Hạn dùng: 10-25 năm |
                           +------------------------+
                                        |
                         Ký cấp phát    | 1-2 lần trong đời
                                        v
                 +--------------------------------------------+
                 |          INTERMEDIATE ISSUING CA           |
                 |  - Online 24/7 (ACME / Web Enrollment)     |
                 |  - Nằm sau tường lửa DMZ, bảo vệ bởi HSM   |
                 |  - Hạn dùng: 2-5 năm                       |
                 +--------------------------------------------+
                          /                           \
         Ký cấp hàng ngày /                             \ Ký cấp hàng ngày
                         v                               v
       +-------------------------------+   +-------------------------------+
       |       TLS Web Server          |   |      mTLS Client / User       |
       |  - SAN: *.kma.edu.vn          |   |  - Email: hocvien@kma.edu.vn  |
       |  - Hạn dùng: 90 ngày (hoặc 1y)|   |  - Hạn dùng: 1 năm            |
       +-------------------------------+   +-------------------------------+
\end{verbatim}

\subsubsection{So Sánh Các Mô Hình Kiến
Trúc:}\label{so-suxe1nh-cuxe1c-muxf4-huxecnh-kiux1ebfn-truxfac}

\begin{enumerate}
\def\labelenumi{\arabic{enumi}.}
\tightlist
\item
  \textbf{Kiến trúc Single-Tier (Đơn cấp - Rủi ro đặc biệt nghiêm
  trọng):} Root CA ký trực tiếp cho tất cả máy chủ. Nếu máy chủ Root CA
  bị xâm nhập hoặc lộ khóa, toàn bộ hạ tầng sụp đổ và không thể thu hồi
  mà phải cài đặt lại Trust Store trên toàn bộ thiết bị người dùng.
\item
  \textbf{Kiến trúc Two-Tier (Hai cấp - Doanh nghiệp vừa):} 1 Root CA
  offline và 1 Issuing Intermediate CA online. Khi Issuing CA gặp sự cố,
  Root CA phát hành một bản thu hồi CRL và ký cấp một Issuing CA mới mà
  không làm gián đoạn Trust Anchor của người dùng.
\item
  \textbf{Kiến trúc Three-Tier (Ba cấp - Chuẩn Quốc gia \& Tài chính):}
  Root CA \(\to\) Policy CA (Phân tách theo địa lý hoặc phân cấp bảo
  mật) \(\to\) Issuing CA \(\to\) End-Entity. Tách biệt hoàn toàn quyền
  hạn vận hành và kiểm toán.
\end{enumerate}

\begin{center}\rule{0.5\linewidth}{0.5pt}\end{center}

\subsection{3.2. Nghi Thức Lễ Sinh Khóa (Key Ceremony) \& Thiết Bị Phần
Cứng
HSM}\label{nghi-thux1ee9c-lux1ec5-sinh-khuxf3a-key-ceremony-thiux1ebft-bux1ecb-phux1ea7n-cux1ee9ng-hsm}

Khóa bí mật của Root CA là linh hồn của hệ thống PKI. Việc tạo ra cặp
khóa này không thể thực hiện tùy tiện trên máy tính cá nhân mà phải
thông qua một nghi thức pháp lý và kỹ thuật trang trọng gọi là
\textbf{Key Ceremony}.

\subsubsection{3.2.1. Quy Trình 6 Bước Của Một Buổi Lễ Key Ceremony
Chuẩn Quốc
Tế}\label{quy-truxecnh-6-bux1b0ux1edbc-cux1ee7a-mux1ed9t-buux1ed5i-lux1ec5-key-ceremony-chuux1ea9n-quux1ed1c-tux1ebf}

\begin{enumerate}
\def\labelenumi{\arabic{enumi}.}
\tightlist
\item
  \textbf{Chuẩn bị Phòng Kín Cách Ly (Faraday Clean Room):} Phòng tổ
  chức phải được quét sóng điện từ, không có thiết bị thu phát vô tuyến,
  cách ly mạng hoàn toàn (Air-gapped), có hệ thống camera giám sát góc
  rộng ghi hình liên tục không góc chết.
\item
  \textbf{Kiểm Tra Niêm Phong Phần Cứng:} Các thiết bị HSM mới nguyên
  đai nguyên kiện được kiểm toán viên độc lập (Independent Auditor) kiểm
  tra số niêm phong (tamper-evident seal) và chứng thực chưa từng bị can
  thiệp.
\item
  \textbf{Khởi Động Môi Trường Live OS Sạch:} Máy chủ điều khiển nạp hệ
  điều hành chuyên dụng từ đĩa quang DVD-R (Write-Once) hoặc thẻ USB chỉ
  đọc có chữ ký số xác thực tính toàn vẹn của mã nhị phân.
\item
  \textbf{Phân Quyền Kiểm Soát Kép (Dual Control / M-of-N):} Khởi tạo
  thẻ quản trị viên HSM theo thuật toán phân chia bí mật Shamir: Phải có
  ít nhất \(M\) trên tổng số \(N\) người giữ khóa (Key Custodians) cùng
  cắm thẻ thông minh và nhập mã PIN thì mới mở được phân vùng bảo mật.
\item
  \textbf{Sinh Cặp Khóa Phần Cứng:} HSM tự kích hoạt bộ tạo số ngẫu
  nhiên vật lý (Hardware TRNG) để sinh cặp khóa Root bên trong bộ nhớ
  được bọc bảo vệ. Khóa bí mật \textbf{TUYỆT ĐỐI KHÔNG BAO GIỜ RA KHỎI
  HSM DƯỚI DẠNG BẢN RÕ (Plaintext)}.
\item
  \textbf{Ký Tự Chứng Thực \& Niêm Phong Két Sắt:} Xuất chứng chỉ Root
  CA công khai, sao lưu bản mã hóa khóa sang thẻ bảo mật dự phòng, sau
  đó ngắt điện HSM, đóng gói niêm phong và lưu trữ trong hai két sắt
  chống cháy tại hai địa điểm địa lý khác nhau.
\end{enumerate}

\subsubsection{3.2.2. Thiết Bị Bảo Mật Phần Cứng HSM \& Chuẩn Giao Tiếp
PKCS\#11}\label{thiux1ebft-bux1ecb-bux1ea3o-mux1eadt-phux1ea7n-cux1ee9ng-hsm-chuux1ea9n-giao-tiux1ebfp-pkcs11}

\begin{itemize}
\item
  \textbf{Tiêu chuẩn NIST FIPS 140-2 / FIPS 140-3:}

  \begin{itemize}
  \tightlist
  \item
    \emph{Level 1:} Bảo vệ phần mềm cơ bản.
  \item
    \emph{Level 2:} Vỏ kim loại có khóa và tem niêm phong chống cạy mở.
  \item
    \emph{Level 3:} Kháng can thiệp tích cực --- Tự động hủy khóa ngay
    lập tức (Zeroization) nếu phát hiện hành vi cạy nắp hoặc tháo vỏ bảo
    vệ.
  \item
    \emph{Level 4:} Bảo vệ mức vật lý tối thượng --- Tự hủy khóa nếu
    phát hiện thay đổi nhiệt độ đột ngột (tấn công làm lạnh sâu) hoặc
    biến đổi điện áp nguồn.
  \end{itemize}
\item
  \textbf{Chuẩn giao diện PKCS\#11 (Cryptoki):} Định nghĩa API chuẩn C
  độc lập với nhà sản xuất (Thales, Entrust, Yubico) cho phép các ứng
  dụng (OpenSSL, Step-CA, HashiCorp Vault) ra lệnh cho HSM ký số mà
  không bao giờ cần đọc khóa bí mật vào RAM máy tính:

\begin{Shaded}
\begin{Highlighting}[]
\NormalTok{C\_Initialize}\OperatorTok{(...);}
\NormalTok{C\_OpenSession}\OperatorTok{(}\NormalTok{slotID}\OperatorTok{,}\NormalTok{ CKF\_SERIAL\_SESSION }\OperatorTok{|}\NormalTok{ CKF\_RW\_SESSION}\OperatorTok{,} \OperatorTok{...);}
\NormalTok{C\_Login}\OperatorTok{(}\NormalTok{hSession}\OperatorTok{,}\NormalTok{ CKU\_USER}\OperatorTok{,}\NormalTok{ pin}\OperatorTok{,}\NormalTok{ pinLen}\OperatorTok{);}
\NormalTok{C\_SignInit}\OperatorTok{(}\NormalTok{hSession}\OperatorTok{,} \OperatorTok{\&}\NormalTok{mechanism}\OperatorTok{,}\NormalTok{ hPrivateKey}\OperatorTok{);}
\NormalTok{C\_Sign}\OperatorTok{(}\NormalTok{hSession}\OperatorTok{,}\NormalTok{ pData}\OperatorTok{,}\NormalTok{ ulDataLen}\OperatorTok{,}\NormalTok{ pSignature}\OperatorTok{,} \OperatorTok{\&}\NormalTok{ulSignatureLen}\OperatorTok{);}
\end{Highlighting}
\end{Shaded}
\end{itemize}

\begin{center}\rule{0.5\linewidth}{0.5pt}\end{center}

\subsection{3.3. Khung Quy Chế Chứng Thực CP \& CPS (RFC
3647)}\label{khung-quy-chux1ebf-chux1ee9ng-thux1ef1c-cp-cps-rfc-3647}

Một hệ thống PKI chỉ có giá trị pháp lý khi đi kèm với hai văn kiện nền
tảng: * \textbf{Chính Sách Chứng Thực (Certificate Policy - CP):} Tuyên
ngôn cấp quản lý nêu rõ chứng chỉ được cấp cho ai, áp dụng cho mục đích
gì và mức độ tin cậy đến đâu. * \textbf{Quy Chế Thực Hành Chứng Thực
(Certification Practice Statement - CPS):} Tài liệu hướng dẫn kỹ thuật
chi tiết tường tận cách thức CA vận hành để đạt được các mục tiêu đề ra
trong CP.

RFC 3647 chuẩn hóa cấu trúc của CP và CPS thành \textbf{8 đề mục cốt
lõi}: 1. \textbf{Giới thiệu chung (Introduction):} Nhận diện tài liệu,
các bên tham gia (CA, RA, Thuê bao Subscriber, Người tin cậy Relying
Party), phạm vi khả dụng. 2. \textbf{Trách nhiệm công bố và lưu trữ
(Publication \& Repository Controls):} Địa chỉ kho lưu trữ trực tuyến,
tần suất cập nhật CRL/CPS. 3. \textbf{Định danh và xác thực
(Identification \& Authentication):} Thủ tục xác minh danh tính tổ
chức/cá nhân khi đăng ký lần đầu và khi thu hồi. 4. \textbf{Vòng đời
hoạt động của chứng chỉ (Certificate Life-Cycle Operations):} Quy trình
nộp đơn, cấp phát, gia hạn và thu hồi chứng chỉ. 5. \textbf{Kiểm soát cơ
sở vật chất, quản lý và vận hành (Facility, Management Controls):} An
ninh phòng máy chủ, kiểm soát truy cập vật lý, phòng cháy chữa cháy, sao
lưu thảm họa. 6. \textbf{Kiểm soát kỹ thuật an ninh (Technical Security
Controls):} Chuẩn thế hệ khóa, độ dài khóa, chuẩn HSM, cơ chế kiểm soát
kép, an ninh mạng phân vùng DMZ. 7. \textbf{Hồ sơ định dạng
(Certificate, CRL, and OCSP Profiles):} Cấu trúc chính xác các OID,
thuật toán ký, các trường mở rộng Critical/Non-critical. 8. \textbf{Đánh
giá và kiểm toán tuân thủ (Compliance Audit \& Assessments):} Cơ quan
kiểm toán độc lập, chu kỳ kiểm toán hàng năm (WebTrust, ETSI EN 319
411).

\begin{center}\rule{0.5\linewidth}{0.5pt}\end{center}

\section{CHƯƠNG 4: THUẬT TOÁN XÁC THỰC ĐƯỜNG DẪN TIN CẬY \& CƠ CHẾ THU
HỒI CHỨNG
CHỈ}\label{chux1b0ux1a1ng-4-thuux1eadt-touxe1n-xuxe1c-thux1ef1c-ux111ux1b0ux1eddng-dux1eabn-tin-cux1eady-cux1a1-chux1ebf-thu-hux1ed3i-chux1ee9ng-chux1ec9}

\subsection{4.1. Thuật Toán Xác Thực Chuỗi Chứng Chỉ (RFC 5280
§6)}\label{thuux1eadt-touxe1n-xuxe1c-thux1ef1c-chuux1ed7i-chux1ee9ng-chux1ec9-rfc-5280-6}

Khi người dùng truy cập một website bảo mật, trình duyệt phải dựng một
chuỗi chứng chỉ từ chứng chỉ lá \(C_0\) ngược về một chứng chỉ gốc
\(C_n\) nằm trong cơ sở dữ liệu tin cậy cục bộ (Local Trust Store):

\[C_0 (\text{Leaf}) \to C_1 (\text{Issuing CA}) \to C_2 (\text{Intermediate CA}) \to C_n (\text{Trust Anchor / Root CA})\]

Thuật toán xác thực đường dẫn (Path Validation Algorithm) duyệt qua từng
mắt xích từ \(i = 1\) đến \(n\) và thực hiện tuần tự \textbf{6 bước kiểm
định toán học nghiêm ngặt}: 1. \textbf{Kiểm tra chữ ký số:} Sử dụng khóa
công khai của chứng chỉ cha \(C_i\) để giải mã và kiểm tra chữ ký trên
khối TBSCertificate của chứng chỉ con \(C_{i-1}\). 2. \textbf{Kiểm tra
thời gian hiệu lực:} Thời điểm hiện tại \(T_{\text{now}}\) phải nằm
trong khoảng:
\[C_{i-1}.\text{notBefore} \le T_{\text{now}} \le C_{i-1}.\text{notAfter}\]
3. \textbf{Khớp nối chuỗi định danh (Name Chaining):} Tên định danh của
người phát hành trong chứng chỉ con phải trùng khớp hoàn toàn từng byte
nhị phân với tên chủ thể của chứng chỉ cha:
\[C_{i-1}.\text{Issuer} == C_i.\text{Subject}\] 4. \textbf{Kiểm tra
quyền hạn CA (Basic Constraints):} Đối với tất cả các mắt xích trung
gian (\(i \ge 1\)), chứng chỉ \(C_i\) bắt buộc phải có phần mở rộng
\texttt{Basic\ Constraints} với cờ \texttt{cA\ =\ TRUE}. 5. \textbf{Kiểm
tra giới hạn độ sâu (Path Length Constraint):} Nếu \(C_i\) có quy định
\texttt{pathLenConstraint\ =\ k}, số lượng chứng chỉ CA trung gian nằm
giữa nó và chứng chỉ thực thể cuối không được vượt quá \(k\). 6.
\textbf{Kiểm tra trạng thái thu hồi (Revocation Status):} Kiểm tra xem
số seri của \(C_{i-1}\) có bị liệt kê trong danh sách thu hồi CRL hoặc
nhận phản hồi \texttt{REVOKED} từ OCSP Responder hay không.

\begin{center}\rule{0.5\linewidth}{0.5pt}\end{center}

\subsection{4.2. Cơ Chế Thu Hồi: CRL vs.~OCSP vs.~OCSP
Stapling}\label{cux1a1-chux1ebf-thu-hux1ed3i-crl-vs.-ocsp-vs.-ocsp-stapling}

Một chứng chỉ có thể bị hủy bỏ giá trị trước khi hết hạn do nhiều nguyên
nhân: lộ khóa bí mật (\texttt{keyCompromise}), thay đổi thông tin định
danh (\texttt{affiliationChanged}), hoặc đơn vị ngừng hoạt động. Ba cơ
chế kiểm tra thu hồi được phát triển qua các thời kỳ:

\begin{verbatim}
1. CRL (Danh sách thu hồi tĩnh):
   Client -------- Tải toàn bộ file CRL (dung lượng lớn hàng MB) --------> CA CRL Server
   (Nhược điểm: Tải nặng, độ trễ cập nhật lớn, gây chậm kết nối)

2. OCSP (Tra cứu trực tuyến):
   Client -------- Truy vấn: Chứng chỉ Serial #12345 còn sống không? -------> OCSP Responder
          <------- Phản hồi: Trạng thái GOOD (ký bởi CA) ------------------
   (Nhược điểm: Lộ lịch sử duyệt web của Client cho CA; nguy cơ Soft-Fail bị chặn)

3. OCSP STAPLING (Khuyến nghị chuẩn hiện đại TLS 1.3):
   Web Server tự truy vấn OCSP Responder định kỳ (mỗi giờ) và lưu đệm phản hồi.
   Client <======= Gói bắt tay TLS 1.3 kèm luôn phản hồi OCSP đã ký ======== Web Server
   (Ưu điểm tuyệt đối: Tốc độ tối đa 0-RTT, bảo mật riêng tư 100%, không nghẽn CA)
\end{verbatim}

\subsubsection{So Sánh Kỹ Thuật Chi
Tiết:}\label{so-suxe1nh-kux1ef9-thuux1eadt-chi-tiux1ebft}

{\def\LTcaptype{none} % do not increment counter
\begin{longtable}[]{@{}
  >{\raggedright\arraybackslash}p{(\linewidth - 6\tabcolsep) * \real{0.2500}}
  >{\raggedright\arraybackslash}p{(\linewidth - 6\tabcolsep) * \real{0.2500}}
  >{\raggedright\arraybackslash}p{(\linewidth - 6\tabcolsep) * \real{0.2500}}
  >{\raggedright\arraybackslash}p{(\linewidth - 6\tabcolsep) * \real{0.2500}}@{}}
\toprule\noalign{}
\begin{minipage}[b]{\linewidth}\raggedright
Tiêu Chí Kỹ Thuật
\end{minipage} & \begin{minipage}[b]{\linewidth}\raggedright
CRL (RFC 5280)
\end{minipage} & \begin{minipage}[b]{\linewidth}\raggedright
OCSP (RFC 6960)
\end{minipage} & \begin{minipage}[b]{\linewidth}\raggedright
OCSP Stapling (RFC 6066 / 7633)
\end{minipage} \\
\midrule\noalign{}
\endhead
\bottomrule\noalign{}
\endlastfoot
\textbf{Phương thức truyền tải} & File tĩnh qua HTTP/FTP & Giao thức nhị
phân HTTP POST/GET & Kẹp trong gói TLS Handshake \\
\textbf{Băng thông Client} & Rất nặng (tải danh sách lớn) & Rất nhẹ (vài
trăm bytes) & Không tốn thêm kết nối ngoài \\
\textbf{Quyền riêng tư (Privacy)} & Tốt (CA không biết ai bị tra cứu) &
Kém (CA biết IP người dùng truy cập web nào) & Hoàn hảo (CA không biết
danh tính client) \\
\textbf{Độ trễ kết nối mạng} & Lớn (nếu file CRL chưa cache) & Thêm 1
RTT tới OCSP server & 0 RTT (có sẵn trong gói TLS) \\
\textbf{Hành vi khi mạng lỗi} & Soft-fail (bỏ qua kiểm tra) & Soft-fail
(nguy cơ MitM chặn) & \textbf{Hard-fail} khi bật Must-Staple \\
\end{longtable}
}

\begin{center}\rule{0.5\linewidth}{0.5pt}\end{center}

\section{CHƯƠNG 5: GIAO THỨC QUẢN TRỊ TỰ ĐỘNG ACME (RFC 8555) \& HẠ TẦNG
ZERO TRUST
PKI}\label{chux1b0ux1a1ng-5-giao-thux1ee9c-quux1ea3n-trux1ecb-tux1ef1-ux111ux1ed9ng-acme-rfc-8555-hux1ea1-tux1ea7ng-zero-trust-pki}

\subsection{5.1. Cuộc Cách Mạng Tự Động Hóa Vòng Đời Chứng Chỉ Với ACME
(RFC
8555)}\label{cuux1ed9c-cuxe1ch-mux1ea1ng-tux1ef1-ux111ux1ed9ng-huxf3a-vuxf2ng-ux111ux1eddi-chux1ee9ng-chux1ec9-vux1edbi-acme-rfc-8555}

Năm 2020, Apple, Google và Mozilla chính thức áp đặt giới hạn thời hạn
sống của chứng chỉ TLS xuống còn \textbf{398 ngày}. Năm 2024, Google đề
xuất lộ trình rút ngắn tiếp xuống \textbf{90 ngày}, và trong tương lai
là \textbf{45 ngày}. Việc quản lý chứng chỉ thủ công bằng con người (tạo
CSR, nộp đơn, chờ duyệt, tải về, cấu hình máy chủ) hoàn toàn phá sản.

Giao thức \textbf{ACME (Automated Certificate Management Environment -
RFC 8555)} do Let's Encrypt tiên phong và IETF chuẩn hóa đã giải quyết
bài toán này thông qua tự động hóa hoàn toàn 100\%:

\begin{verbatim}
[ACME Client: Certbot / Step]                      [ACME Server: Step-CA / Let's Encrypt]
             |                                                          |
             | 1. Khởi tạo tài khoản (Account Key Pair Ed25519)         |
             |--------------------------------------------------------->|
             |                                                          |
             | 2. Yêu cầu cấp chứng chỉ cho "api.kma.edu.vn"            |
             |--------------------------------------------------------->|
             | 3. Trả về Thử thách (Challenge: Token ngẫu nhiên)       |
             |<---------------------------------------------------------|
             |                                                          |
             | [Client chuẩn bị đáp ứng thử thách]                     |
             | - HTTP-01: Ghi file vào /.well-known/acme-challenge/    |
             | - DNS-01 : Tạo bản ghi TXT _acme-challenge.domain        |
             |                                                          |
             | 4. Báo hiệu đã sẵn sàng kiểm tra                         |
             |--------------------------------------------------------->|
             |                                                          |
             |           5. Server tự kết nối Internet kiểm tra thử thách
             |              (Đúng token và chữ ký tài khoản)           |
             |                                                          |
             | 6. Xác thực thành công -> Cho phép nộp CSR               |
             |<---------------------------------------------------------|
             | 7. Nộp CSR kèm Public Key máy chủ                        |
             |--------------------------------------------------------->|
             | 8. Ký và trả về chứng chỉ hoàn chỉnh X.509               |
             |<---------------------------------------------------------|
\end{verbatim}

\subsubsection{Các Loại Thử Thách ACME Phổ
Biến:}\label{cuxe1c-loux1ea1i-thux1eed-thuxe1ch-acme-phux1ed5-biux1ebfn}

\begin{enumerate}
\def\labelenumi{\arabic{enumi}.}
\tightlist
\item
  \textbf{HTTP-01 Challenge:} Máy chủ ACME truy vấn qua cổng HTTP 80 để
  đọc token. Phù hợp cho máy chủ đơn lẻ, nhưng không cấp được chứng chỉ
  Wildcard (\texttt{*.domain.com}).
\item
  \textbf{DNS-01 Challenge:} Client sử dụng API của máy chủ DNS để tự
  động tạo bản ghi TXT. Đây là phương pháp mạnh mẽ nhất, hỗ trợ chứng
  chỉ Wildcard và hoạt động tốt ngay cả với máy chủ nằm sâu trong mạng
  nội bộ không mở cổng ra Internet.
\item
  \textbf{TLS-ALPN-01 Challenge:} Dành riêng cho các hệ thống reverse
  proxy (HAProxy, Traefik), đàm phán qua giao thức TLS ALPN chuyên dụng
  trên cổng 443 mà không cần can thiệp tầng HTTP.
\end{enumerate}

\begin{center}\rule{0.5\linewidth}{0.5pt}\end{center}

\subsection{5.2. PKI Trong Kiến Trúc Mạng Không Tin Cậy (Zero Trust
PKI)}\label{pki-trong-kiux1ebfn-truxfac-mux1ea1ng-khuxf4ng-tin-cux1eady-zero-trust-pki}

Mô hình phòng thủ chu vi truyền thống (``bên ngoài nguy hiểm, bên trong
an toàn'') đã hoàn toàn lỗi thời. Kiến trúc \textbf{Zero Trust} áp dụng
nguyên tắc: \textbf{Luôn luôn xác thực, không bao giờ tin tưởng bất kỳ
ai dù đang ở trong mạng nội bộ}.

\begin{itemize}
\tightlist
\item
  \textbf{mTLS (Mutual TLS) Toàn Diện:} Mọi kết nối giữa microservices A
  và microservices B đều bắt buộc cả hai bên phải xuất trình chứng chỉ
  X.509 để xác thực hai chiều và mã hóa đường truyền.
\item
  \textbf{Chứng Chỉ Biểu Kiến (Ephemeral Certificates):} Khóa bí mật và
  chứng chỉ được sinh tự động bởi HashiCorp Vault hoặc Step-CA với thời
  gian sống cực ngắn (từ 1 giờ đến 24 giờ). Khi chứng chỉ hết hạn, hệ
  thống tự động xoay vòng (key rotation). Kẻ tấn công nếu có đánh cắp
  được khóa bí mật trong bộ nhớ thì giá trị sử dụng cũng đã hết hạn sau
  ít phút.
\end{itemize}

\begin{center}\rule{0.5\linewidth}{0.5pt}\end{center}

\section{CHƯƠNG 6: CÁC KỊCH BẢN TẤN CÔNG HẠ TẦNG PKI \& CHIẾN LƯỢC PHÒNG
THỦ TOÀN
DIỆN}\label{chux1b0ux1a1ng-6-cuxe1c-kux1ecbch-bux1ea3n-tux1ea5n-cuxf4ng-hux1ea1-tux1ea7ng-pki-chiux1ebfn-lux1b0ux1ee3c-phuxf2ng-thux1ee7-touxe0n-diux1ec7n}

\subsection{6.1. Bài Học Xương Máu Từ Các Thảm Họa PKI Lịch
Sử}\label{buxe0i-hux1ecdc-xux1b0ux1a1ng-muxe1u-tux1eeb-cuxe1c-thux1ea3m-hux1ecda-pki-lux1ecbch-sux1eed}

PKI được xây dựng dựa trên sự ủy thác. Nếu một CA gốc trong Trust Store
bị xâm nhập, kẻ tấn công có quyền lực tối thượng để giả mạo bất kỳ dịch
vụ nào trên thế giới.

\begin{enumerate}
\def\labelenumi{\arabic{enumi}.}
\tightlist
\item
  \textbf{Vụ Án DigiNotar (Hà Lan - 2011):}

  \begin{itemize}
  \tightlist
  \item
    \emph{Bối cảnh:} DigiNotar là CA công cộng uy tín của chính phủ Hà
    Lan, nằm trong Trust Store mặc định của Windows, iOS, Android,
    Firefox.
  \item
    \emph{Kịch bản tấn công:} Hacker thâm nhập vào mạng nội bộ của
    DigiNotar do hệ thống thiếu phân tách mạng và mật khẩu quản trị yếu.
    Root CA bị để ở trạng thái trực tuyến. Kẻ tấn công phát hành hơn 500
    chứng chỉ giả mạo cho các tên miền trọng yếu: \texttt{*.google.com},
    \texttt{*.torproject.org}, \texttt{*.microsoft.com}.
  \item
    \emph{Hậu quả:} Hơn 300.000 người dùng Internet tại Iran đã bị tấn
    công Man-in-the-Middle đọc trộm toàn bộ email và tài khoản cá nhân.
    Google và Microsoft đã tung bản vá khẩn cấp loại bỏ hoàn toàn
    DigiNotar khỏi Trust Store, dẫn tới việc công ty này nộp đơn phá sản
    sau đó 1 tháng.
  \end{itemize}
\item
  \textbf{Vũ Khí Mạng Flame (2012):}

  \begin{itemize}
  \tightlist
  \item
    \emph{Kịch bản:} Nhóm phát triển mã độc gián điệp Flame đã thực hiện
    một cuộc tấn công toán học kỳ công: Khai thác lỗ hổng va chạm tiền
    ảnh (chosen-prefix collision) trên hàm băm MD5 mà Microsoft vẫn sử
    dụng để ký chứng chỉ dịch vụ Terminal Server.
  \item
    \emph{Khai thác:} Kẻ tấn công tạo ra một chứng chỉ giả mạo có cùng
    mã băm MD5 với chứng chỉ hợp lệ của Microsoft, từ đó ký mã độc giả
    dạng bản cập nhật chính thức của Windows Update để lây nhiễm vào các
    cơ sở hạ tầng trọng yếu tại Trung Đông.
  \end{itemize}
\end{enumerate}

\begin{center}\rule{0.5\linewidth}{0.5pt}\end{center}

\subsection{6.2. Kỹ Thuật Tấn Công Active Directory Certificate Services
(AD
CS)}\label{kux1ef9-thuux1eadt-tux1ea5n-cuxf4ng-active-directory-certificate-services-ad-cs}

Active Directory Certificate Services (AD CS) là thành phần PKI phổ biến
nhất trong các tập đoàn lớn toàn cầu. Năm 2021, Will Schroeder và Lee
Christensen (SpecterOps) đã công bố nghiên cứu gây chấn động chỉ ra hàng
loạt cấu hình sai sót nghiêm trọng cho phép leo thang đặc quyền từ người
dùng cấp thấp lên Quản trị viên miền (Domain Admin):

\begin{itemize}
\tightlist
\item
  \textbf{Kỹ thuật ESC1 (Misconfigured Certificate Template):}

  \begin{itemize}
  \tightlist
  \item
    Mẫu chứng chỉ (Template) thỏa mãn đồng thời 3 điều kiện chết người:

    \begin{enumerate}
    \def\labelenumi{\arabic{enumi}.}
    \tightlist
    \item
      Cấp quyền cho nhóm người dùng bình thường (\texttt{Domain\ Users})
      được nộp đơn xin cấp chứng chỉ.
    \item
      Bật cờ \texttt{ENROLLEE\_SUPPLIES\_SUBJECT}: Cho phép người nộp
      đơn tùy ý điền tên Subject Alternative Name (SAN) trong CSR.
    \item
      Cấu hình Extended Key Usage (EKU) có chứa
      \texttt{Client\ Authentication} (\texttt{1.3.6.1.5.5.7.3.2}) hoặc
      \texttt{Smart\ Card\ Logon}.
    \end{enumerate}
  \item
    \emph{Khai thác:} Kẻ tấn công sử dụng công cụ \texttt{Certipy} yêu
    cầu một chứng chỉ với SAN giả mạo là tài khoản
    \texttt{Administrator@domain.local}. Sau khi nhận được chứng chỉ từ
    AD CS, kẻ tấn công dùng chứng chỉ đó yêu cầu vé Kerberos TGT (Ticket
    Granting Ticket) để chiếm quyền điều khiển toàn bộ mạng công ty.
  \end{itemize}
\end{itemize}

\begin{Shaded}
\begin{Highlighting}[]
\CommentTok{\# Lệnh trinh sát và khai thác ESC1 bằng công cụ Certipy:}
\ExtensionTok{certipy}\NormalTok{ find }\AttributeTok{{-}u}\NormalTok{ user@corp.local }\AttributeTok{{-}p}\NormalTok{ Password123 }\AttributeTok{{-}dc{-}ip}\NormalTok{ 10.0.0.1 }\AttributeTok{{-}vulnerable}
\ExtensionTok{certipy}\NormalTok{ req }\AttributeTok{{-}u}\NormalTok{ user@corp.local }\AttributeTok{{-}p}\NormalTok{ Password123 }\AttributeTok{{-}ca}\NormalTok{ Corp{-}CA }\AttributeTok{{-}template}\NormalTok{ ESC1{-}Template }\AttributeTok{{-}upn}\NormalTok{ administrator@corp.local}
\ExtensionTok{certipy}\NormalTok{ auth }\AttributeTok{{-}pfx}\NormalTok{ administrator.pfx }\AttributeTok{{-}dc{-}ip}\NormalTok{ 10.0.0.1}
\end{Highlighting}
\end{Shaded}

\begin{itemize}
\tightlist
\item
  \textbf{Biện Pháp Phòng Thủ:}

  \begin{itemize}
  \tightlist
  \item
    Bỏ cờ \texttt{ENROLLEE\_SUPPLIES\_SUBJECT} trên mọi template có tính
    năng Client Authentication.
  \item
    Thiết lập cơ chế phê duyệt thủ công
    (\texttt{CA\ administrator\ approval}) cho các chứng chỉ nhạy cảm.
  \item
    Giám sát Windows Event Log ID \texttt{4886} (Yêu cầu cấp chứng chỉ)
    và \texttt{4887} (Đã cấp chứng chỉ) bằng các luật phát hiện
    \textbf{Sigma rules}.
  \end{itemize}
\end{itemize}

\begin{center}\rule{0.5\linewidth}{0.5pt}\end{center}

\subsection{6.3. Hệ Thống Minh Bạch Chứng Chỉ (Certificate Transparency
- RFC
6962)}\label{hux1ec7-thux1ed1ng-minh-bux1ea1ch-chux1ee9ng-chux1ec9-certificate-transparency---rfc-6962}

Để giải quyết triệt để vấn đề CA cấp lén lút chứng chỉ giả mạo như vụ
DigiNotar, Google đã khởi xướng chuẩn \textbf{Certificate Transparency
(CT Logs - RFC 6962)}:

\begin{enumerate}
\def\labelenumi{\arabic{enumi}.}
\tightlist
\item
  \textbf{Nguyên Lý Hoạt Động:} Mọi chứng chỉ muốn được trình duyệt
  (Chrome, Safari) tin cậy \textbf{BẮT BUỘC} phải được ghi vào các máy
  chủ CT Log công khai trước khi phát hành.
\item
  \textbf{Cấu Trúc Cây Merkle (Merkle Hash Tree):}

  \begin{itemize}
  \tightlist
  \item
    CT Log là một cuốn sổ cái chỉ thêm (Append-Only Cryptographic
    Ledger). Dữ liệu được tổ chức dưới dạng cây nhị phân Merkle Tree.
  \item
    Mỗi nút lá là một chứng chỉ đã được nộp. Nút cha là mã băm của hai
    nút con:
    \(N_{\text{parent}} = H(N_{\text{left}} \parallel N_{\text{right}})\).
  \item
    Gốc cây (Merkle Tree Root) đại diện cho toàn bộ trạng thái lịch sử
    của hàng tỷ chứng chỉ đã phát hành.
  \end{itemize}
\end{enumerate}

\begin{verbatim}
                         Merkle Root Hash (Ký bởi CT Log Server)
                                    /              \
                                   /                \
                             Node 12                Node 34
                             /      \               /      \
                          Hash 1   Hash 2        Hash 3   Hash 4
                           [C1]     [C2]          [C3]     [C4]
\end{verbatim}

\begin{enumerate}
\def\labelenumi{\arabic{enumi}.}
\setcounter{enumi}{2}
\tightlist
\item
  \textbf{Signed Certificate Timestamp (SCT):} Khi CA gửi chứng chỉ tạm
  (Pre-certificate) lên CT Log, Log server trả về một biên nhận mật mã
  gọi là SCT. SCT là bằng chứng cam kết rằng chứng chỉ này sẽ được đưa
  vào cây Merkle trong vòng tối đa 24 giờ (Maximum Merge Delay - MMD).
  SCT được nhúng trực tiếp vào chứng chỉ X.509 chính thức dưới dạng một
  Extension.
\item
  \textbf{Kiểm Định Tuân Thủ Với ZLint:} Trước khi đưa chứng chỉ ra môi
  trường production, các tổ chức CA sử dụng công cụ nguồn mở
  \textbf{ZLint} để quét tự động hơn 250 điều khoản kỹ thuật của chuẩn
  CA/B Forum, bảo đảm 100\% chứng chỉ không mắc các lỗi sai sót về định
  dạng bit hay OID.
\end{enumerate}

\begin{center}\rule{0.5\linewidth}{0.5pt}\end{center}

\section{CHƯƠNG 7: CHUYỂN DỊCH HỆ THỐNG SANG MẬT MÃ HẬU LƯỢNG TỬ (PQC
MIGRATION)}\label{chux1b0ux1a1ng-7-chuyux1ec3n-dux1ecbch-hux1ec7-thux1ed1ng-sang-mux1eadt-muxe3-hux1eadu-lux1b0ux1ee3ng-tux1eed-pqc-migration}

\subsection{7.1. Mối Đe Dọa Từ Máy Tính Lượng Tử \& Chiến Lược
HNDL}\label{mux1ed1i-ux111e-dux1ecda-tux1eeb-muxe1y-tuxednh-lux1b0ux1ee3ng-tux1eed-chiux1ebfn-lux1b0ux1ee3c-hndl}

Năm 1994, Peter Shor công bố thuật toán lượng tử chứng minh rằng: Một
máy tính lượng tử lý tưởng (Cryptanalytically Relevant Quantum Computer
- CRQC) có thể tìm chu kỳ hàm số và giải quyết đồng thời bài toán phân
tích thừa số nguyên tố và logarit rời rạc trong \textbf{thời gian đa
thức}:

\[\mathcal{O}\left((\log N)^3\right)\]

Điều này đồng nghĩa với việc toàn bộ thuật toán RSA, DSA, ECDSA và ECDH
đang bảo vệ Internet \textbf{sẽ hoàn toàn sụp đổ chỉ sau một đêm} khi
máy tính lượng tử đạt đủ số lượng Qubit vật lý ổn định.

Mối đe dọa không nằm ở tương lai 10 năm nữa mà đang diễn ra \textbf{ngay
trong giây phút này} thông qua chiến lược tấn công:
\[\textbf{Harvest Now, Decrypt Later (HNDL) — Thu thập ngay, Giải mã sau}\]
Các cơ quan tình báo và nhóm tội phạm mạng tinh vi đang bí mật thu thập
và lưu trữ hàng terabyte lưu lượng mã hóa HTTPS, dữ liệu ngoại giao, hồ
sơ bệnh án và giao dịch ngân hàng của các quốc gia. Khi máy tính lượng
tử ra đời, toàn bộ kho dữ liệu lịch sử này sẽ bị giải mã hồi tố, làm lộ
toàn bộ bí mật nhà nước và bí mật kinh doanh!

\begin{center}\rule{0.5\linewidth}{0.5pt}\end{center}

\subsection{7.2. Bộ Tiêu Chuẩn Mật Mã Hậu Lượng Tử Chính Thức Của NIST
(08/2024)}\label{bux1ed9-tiuxeau-chuux1ea9n-mux1eadt-muxe3-hux1eadu-lux1b0ux1ee3ng-tux1eed-chuxednh-thux1ee9c-cux1ee7a-nist-082024}

Sau cuộc thi tuyển chọn toàn cầu kéo dài gần một thập kỷ, ngày
13/08/2024, Viện Tiêu chuẩn và Công nghệ Quốc gia Hoa Kỳ (NIST) đã chính
thức ban hành bộ tiêu chuẩn liên bang:

\subsubsection{7.2.1. NIST FIPS 203: Chuẩn Đóng Gói Khóa ML-KEM
(Module-Lattice
KEM)}\label{nist-fips-203-chuux1ea9n-ux111uxf3ng-guxf3i-khuxf3a-ml-kem-module-lattice-kem}

\begin{itemize}
\tightlist
\item
  Dựa trên thuật toán nguyên bản \textbf{CRYSTALS-Kyber}.
\item
  Dựa trên bài toán khó \textbf{Module Learning With Errors (M-LWE)}
  trên vành đa thức hữu hạn: Cho ma trận đa thức \(\mathbf{A}\), vector
  bí mật \(\mathbf{s}\) và vector sai số nhỏ \(\mathbf{e}\), việc tìm
  \(\mathbf{s}\) từ \(\mathbf{b} = \mathbf{A}\mathbf{s} + \mathbf{e}\)
  là bài toán NP-hard ngay cả với máy tính lượng tử.
\item
  \textbf{Ba cấp độ an ninh:}

  \begin{itemize}
  \tightlist
  \item
    \texttt{ML-KEM-512}: Tương đương mức an ninh AES-128 (NIST Level 1).
  \item
    \texttt{ML-KEM-768}: Tương đương mức an ninh AES-192 (NIST Level 3 -
    Khuyến nghị chuẩn chung).
  \item
    \texttt{ML-KEM-1024}: Tương đương mức an ninh AES-256 (NIST Level 5
    - Dành cho hạ tầng tối mật).
  \end{itemize}
\end{itemize}

\subsubsection{7.2.2. NIST FIPS 204: Chuẩn Chữ Ký Số ML-DSA
(Module-Lattice
DSA)}\label{nist-fips-204-chuux1ea9n-chux1eef-kuxfd-sux1ed1-ml-dsa-module-lattice-dsa}

\begin{itemize}
\tightlist
\item
  Dựa trên thuật toán nguyên bản \textbf{CRYSTALS-Dilithium}.
\item
  Áp dụng kiến trúc Fiat-Shamir with Aborts trên vành mạng lưới
  Module-LWE và Module-SIS.
\item
  \textbf{Ba cấp độ an ninh:} \texttt{ML-DSA-44}, \texttt{ML-DSA-65}
  (Khuyến nghị chuẩn), \texttt{ML-DSA-87}.
\end{itemize}

\subsubsection{7.2.3. NIST FIPS 205: Chuẩn Chữ Ký Số Phi Trạng Thái
SLH-DSA
(SPHINCS+)}\label{nist-fips-205-chuux1ea9n-chux1eef-kuxfd-sux1ed1-phi-trux1ea1ng-thuxe1i-slh-dsa-sphincs}

\begin{itemize}
\tightlist
\item
  Dựa trên toán học hàm băm phi cấu trúc (Stateless Hash-Based).
\item
  Ưu điểm: Cơ sở an toàn hoàn toàn độc lập với toán học mạng lưới (được
  xem là tuyến phòng thủ dự phòng an toàn tuyệt đối nếu một ngày toán
  học mạng lưới bị giải mã). Nhược điểm: Chữ ký có kích thước rất lớn
  (hơn 17 KB - 40 KB).
\end{itemize}

\begin{center}\rule{0.5\linewidth}{0.5pt}\end{center}

\subsection{7.3. Kiến Trúc Bắt Tay Lai (Hybrid TLS 1.3 \& Hybrid
Certificates)}\label{kiux1ebfn-truxfac-bux1eaft-tay-lai-hybrid-tls-1.3-hybrid-certificates}

Việc chuyển dịch từ mật mã cổ điển sang PQC không thể diễn ra đột ngột
mà phải thông qua cơ chế \textbf{Bắt tay lai (Hybrid Approach)} theo
chuẩn IETF:

\begin{verbatim}
[Client]                                                        [Server]
   |                                                               |
   | 1. ClientHello:                                               |
   |    Trao đổi khóa Lai (Supported Groups):                      |
   |    - Khóa cổ điển : X25519 (32 bytes)                         |
   |    - Khóa PQC     : ML-KEM-768 (1184 bytes)                  |
   |-------------------------------------------------------------->|
   |                                                               |
   | 2. ServerHello:                                               |
   |    - Phản hồi X25519 (32 bytes)                               |
   |    - Phản hồi Ciphertext ML-KEM-768 (1088 bytes)              |
   |<--------------------------------------------------------------|
   |                                                               |
   | [Cả hai bên tính toán Khóa bí mật chung kép]:                 |
   | Shared_Secret = KDF(SS_X25519 || SS_ML-KEM-768)               |
   | (Hệ thống an toàn nếu ít nhất một trong hai thuật toán an toàn)|
\end{verbatim}

\subsubsection{Thử Thách Kỹ Thuật Khi Chuyển Dịch Lên Môi Trường Sản
Xuất:}\label{thux1eed-thuxe1ch-kux1ef9-thuux1eadt-khi-chuyux1ec3n-dux1ecbch-luxean-muxf4i-trux1b0ux1eddng-sux1ea3n-xuux1ea5t}

\begin{enumerate}
\def\labelenumi{\arabic{enumi}.}
\tightlist
\item
  \textbf{Bùng Nổ Kích Thước Khóa \& Chữ Ký:}

  \begin{itemize}
  \tightlist
  \item
    Khóa ECDSA P-256: Public Key 64 bytes, Signature 64 bytes.
  \item
    Khóa ML-DSA-65: Public Key \textbf{1952 bytes}, Signature
    \textbf{3309 bytes}.
  \end{itemize}
\item
  \textbf{Hiện Tượng Phân Mảnh Gói Tin TCP (TCP MTU Packet
  Fragmentation):}

  \begin{itemize}
  \tightlist
  \item
    MTU tiêu chuẩn mạng Ethernet là 1500 bytes. Một chứng chỉ hoặc chữ
    ký ML-DSA vượt quá kích thước 1 MTU bắt buộc phải bị phân mảnh thành
    3 - 4 khung tin IP.
  \item
    Trên mạng di động (4G/5G) hoặc kết nối mạng vệ tinh có tỉ lệ mất gói
    cao, việc phân mảnh gói tin làm tăng đột biến độ trễ bắt tay mạng
    (Handshake Latency) và tăng nguy cơ rớt kết nối. Đây là bài toán kỹ
    thuật trọng tâm mà học viên Nhóm 3 sẽ đo kiểm thực tế trong đồ án
    môn học.
  \end{itemize}
\end{enumerate}

\begin{center}\rule{0.5\linewidth}{0.5pt}\end{center}

\section{CHƯƠNG 8: KHUNG PHÁP LÝ, QUY CHUẨN KỸ THUẬT \& HẠ TẦNG CHỮ KÝ
SỐ QUỐC
GIA}\label{chux1b0ux1a1ng-8-khung-phuxe1p-luxfd-quy-chuux1ea9n-kux1ef9-thuux1eadt-hux1ea1-tux1ea7ng-chux1eef-kuxfd-sux1ed1-quux1ed1c-gia}

\subsection{8.1. Hệ Thống Luật Pháp \& Quy Chuẩn Kỹ Thuật Việt
Nam}\label{hux1ec7-thux1ed1ng-luux1eadt-phuxe1p-quy-chuux1ea9n-kux1ef9-thuux1eadt-viux1ec7t-nam}

Tại Việt Nam, chữ ký số và hạ tầng khóa công khai được điều chỉnh bởi
một hành lang pháp lý hoàn chỉnh và chặt chẽ: 1. \textbf{Luật Giao dịch
điện tử số 20/2023/QH15 (Chính thức có hiệu lực từ ngày 01/07/2024):} -
Thay thế toàn diện Luật Giao dịch điện tử năm 2005. - Quy định phân loại
chi tiết: \textbf{Chữ ký điện tử chuyên dùng}, \textbf{Chữ ký số chuyên
dùng công vụ} và \textbf{Chữ ký số công cộng}. - Khẳng định giá trị pháp
lý tương đương chữ ký tay của cá nhân và con dấu của tổ chức. - Bổ sung
quy định về \textbf{Dịch vụ Tin cậy (Trust Services)} bao gồm: Dịch vụ
cấp dấu thời gian (Time-stamping), Dịch vụ chứng thực thông điệp dữ
liệu, và Dịch vụ chứng thực chữ ký số. 2. \textbf{Nghị định
130/2018/NĐ-CP:} Hướng dẫn chi tiết thi hành Luật Giao dịch điện tử về
chữ ký số và dịch vụ chứng thực chữ ký số. 3. \textbf{Thông tư số
16/2019/TT-BTTTT \& Thông tư số 31/2020/TT-BTTTT:} Quy định danh mục
tiêu chuẩn kỹ thuật bắt buộc áp dụng cho chữ ký số và thiết bị lưu khóa:
- Thuật toán băm: Tối thiểu SHA-256 (loại bỏ hoàn toàn SHA-1). - Thuật
toán khóa công khai: RSA tối thiểu 2048-bit hoặc ECC tối thiểu 256-bit.
- Giao thức kiểm tra trạng thái: Bắt buộc hỗ trợ song song cả CRL và
OCSP trực tuyến.

\begin{center}\rule{0.5\linewidth}{0.5pt}\end{center}

\subsection{8.2. Kiến Trúc Hạ Tầng Chữ Ký Số Quốc
Gia}\label{kiux1ebfn-truxfac-hux1ea1-tux1ea7ng-chux1eef-kuxfd-sux1ed1-quux1ed1c-gia}

Hạ tầng PKI Việt Nam được tổ chức theo mô hình hai nhánh quốc gia độc
lập:

\begin{verbatim}
                            +-----------------------------+
                            |       HỆ THỐNG PKI          |
                            |       QUỐC GIA VIỆT NAM     |
                            +-----------------------------+
                                     /             \
                                    /               \
                                   /                 \
  +-----------------------------------+   +-----------------------------------+
  |   TRUNG TÂM CHỨNG THỰC ĐIỆN TỬ    |   |     BAN CƠ YẾU CHÍNH PHỦ          |
  |     QUỐC GIA (NEAC) - BỘ TT&TT    |   |                                   |
  |  - Vận hành National Root CA      |   |  - Quản lý PKI Chuyên Dùng        |
  |  - Cấp phép cho CA Công Cộng      |   |    Công Vụ                        |
  +-----------------------------------+   +-----------------------------------+
                   |                                       |
    Ký cấp phép    |                                       | Ký cấp phát
                   v                                       v
  +-----------------------------------+   +-----------------------------------+
  |      CÁC CA CÔNG CỘNG (PUBLIC CA) |   | CÁC BỘ, BAN, NGÀNH, UBND TỈNH     |
  |  - VNPT-CA, Viettel-CA, FPT-CA,   |   |  - Phục vụ cán bộ, công chức      |
  |    BKAV-CA, TrustCA, MISA-CA...   |   |  - Ký văn bản điều hành điện tử   |
  +-----------------------------------+   +-----------------------------------+
                   |
    Cấp phát cho   |
                   v
  +-----------------------------------+
  |  DOANH NGHIỆP, TỔ CHỨC, CÔNG DÂN  |
  |  - Kê khai thuế, Hải quan điện tử |
  |  - Hóa đơn điện tử, Ngân hàng số  |
  +-----------------------------------+
\end{verbatim}

\begin{center}\rule{0.5\linewidth}{0.5pt}\end{center}

\subsection{8.3. Tiêu Chuẩn Thiết Bị Lưu Khóa Bí Mật \& Ký Số Từ Xa
(Remote
Signing)}\label{tiuxeau-chuux1ea9n-thiux1ebft-bux1ecb-lux1b0u-khuxf3a-buxed-mux1eadt-kuxfd-sux1ed1-tux1eeb-xa-remote-signing}

\begin{enumerate}
\def\labelenumi{\arabic{enumi}.}
\tightlist
\item
  \textbf{Thiết Bị Cứng Cục Bộ (USB Token / Smartcard):}

  \begin{itemize}
  \tightlist
  \item
    Thiết bị phải đạt chứng chỉ an toàn quốc tế \textbf{Common Criteria
    (CC) mức EAL4+} hoặc \textbf{NIST FIPS 140-2 Level 2/3}.
  \item
    Khóa bí mật được sinh trực tiếp bên trong chip bảo mật và không bao
    giờ cho phép sao chép ra ngoài.
  \item
    Có cơ chế khóa tự động sau một số lần nhập sai mã PIN (thường là 5
    lần).
  \end{itemize}
\item
  \textbf{Ký Số Từ Xa (Remote Signing / Cloud HSM):}

  \begin{itemize}
  \tightlist
  \item
    Người dùng không cần mang theo thiết bị phần cứng USB Token cồng
    kềnh. Khóa bí mật được lưu tập trung trong các cụm thiết bị HSM đạt
    chuẩn FIPS 140-2 Level 3 tại trung tâm dữ liệu của nhà cung cấp dịch
    vụ tin cậy (TSP).
  \item
    Xác thực kích hoạt chữ ký thông qua phương thức xác thực đa yếu tố
    (MFA): Sinh trắc học (FIDO2 / FaceID / Vân tay) kết hợp mã OTP thông
    báo qua ứng dụng an toàn.
  \item
    Quy trình ký số từ xa bắt buộc tuân thủ tiêu chuẩn Châu Âu
    \textbf{eIDAS (EN 419 241-1 / EN 419 241-2)} bảo đảm người dùng duy
    trì quyền kiểm soát duy nhất (Sole Control) đối với khóa bí mật của
    mình.
  \end{itemize}
\end{enumerate}

\begin{center}\rule{0.5\linewidth}{0.5pt}\end{center}

\section{PHỤ LỤC A: CẨM NANG 50 LỆNH OPENSSL 3.5 THỰC
CHIẾN}\label{phux1ee5-lux1ee5c-a-cux1ea9m-nang-50-lux1ec7nh-openssl-3.5-thux1ef1c-chiux1ebfn}

\begin{Shaded}
\begin{Highlighting}[]
\CommentTok{\# ==============================================================================}
\CommentTok{\# PHẦN 1: QUẢN LÝ KHÓA CỔ ĐIỂN \& ĐƯỜNG CONG ELLIPTIC}
\CommentTok{\# ==============================================================================}

\CommentTok{\# 1. Sinh khóa bí mật RSA 4096{-}bit mã hóa bằng AES{-}256}
\ExtensionTok{openssl}\NormalTok{ genpkey }\AttributeTok{{-}algorithm}\NormalTok{ RSA }\AttributeTok{{-}pkeyopt}\NormalTok{ rsa\_keygen\_bits:4096 }\AttributeTok{{-}aes{-}256{-}cbc} \AttributeTok{{-}out}\NormalTok{ root\_rsa.key}

\CommentTok{\# 2. Trích xuất khóa công khai từ khóa bí mật RSA}
\ExtensionTok{openssl}\NormalTok{ pkey }\AttributeTok{{-}in}\NormalTok{ root\_rsa.key }\AttributeTok{{-}pubout} \AttributeTok{{-}out}\NormalTok{ root\_rsa.pub}

\CommentTok{\# 3. Liệt kê toàn bộ các đường cong Elliptic được hỗ trợ bởi OpenSSL 3.5}
\ExtensionTok{openssl}\NormalTok{ ecparam }\AttributeTok{{-}list\_curves}

\CommentTok{\# 4. Sinh khóa bí mật ECDSA trên đường cong chuẩn NIST P{-}384}
\ExtensionTok{openssl}\NormalTok{ genpkey }\AttributeTok{{-}algorithm}\NormalTok{ EC }\AttributeTok{{-}pkeyopt}\NormalTok{ ec\_paramgen\_curve:secp384r1 }\AttributeTok{{-}out}\NormalTok{ ca\_ecc.key}

\CommentTok{\# 5. Sinh khóa bí mật Ed25519 (Kháng kênh bên, tối ưu cho TLS 1.3)}
\ExtensionTok{openssl}\NormalTok{ genpkey }\AttributeTok{{-}algorithm}\NormalTok{ ED25519 }\AttributeTok{{-}out}\NormalTok{ server\_ed25519.key}

\CommentTok{\# ==============================================================================}
\CommentTok{\# PHẦN 2: THỰC THI KHỞI TẠO ROOT CA VÀ INTERMEDIATE CA}
\CommentTok{\# ==============================================================================}

\CommentTok{\# 6. Khởi tạo chứng chỉ Root CA tự ký X.509 v3 (Thời hạn 10 năm)}
\ExtensionTok{openssl}\NormalTok{ req }\AttributeTok{{-}x509} \AttributeTok{{-}new} \AttributeTok{{-}nodes} \AttributeTok{{-}key}\NormalTok{ root\_rsa.key }\AttributeTok{{-}sha256} \AttributeTok{{-}days}\NormalTok{ 3650 }\DataTypeTok{\textbackslash{}}
  \AttributeTok{{-}subj} \StringTok{"/C=VN/ST=Ho Chi Minh/O=KMA/OU=Security Division/CN=KMA National Root CA"} \DataTypeTok{\textbackslash{}}
  \AttributeTok{{-}addext} \StringTok{"basicConstraints=critical,CA:TRUE"} \DataTypeTok{\textbackslash{}}
  \AttributeTok{{-}addext} \StringTok{"keyUsage=critical,digitalSignature,keyCertSign,cRLSign"} \DataTypeTok{\textbackslash{}}
  \AttributeTok{{-}addext} \StringTok{"subjectKeyIdentifier=hash"} \DataTypeTok{\textbackslash{}}
  \AttributeTok{{-}out}\NormalTok{ root\_ca.crt}

\CommentTok{\# 7. Tạo yêu cầu ký chứng chỉ (CSR) cho Issuing Intermediate CA}
\ExtensionTok{openssl}\NormalTok{ req }\AttributeTok{{-}new} \AttributeTok{{-}newkey}\NormalTok{ ec }\AttributeTok{{-}pkeyopt}\NormalTok{ ec\_paramgen\_curve:secp384r1 }\AttributeTok{{-}nodes} \DataTypeTok{\textbackslash{}}
  \AttributeTok{{-}keyout}\NormalTok{ intermediate.key }\DataTypeTok{\textbackslash{}}
  \AttributeTok{{-}subj} \StringTok{"/C=VN/ST=Ho Chi Minh/O=KMA/OU=PKI Center/CN=KMA Issuing Intermediate CA"} \DataTypeTok{\textbackslash{}}
  \AttributeTok{{-}out}\NormalTok{ intermediate.csr}

\CommentTok{\# 8. Root CA ký cấp phát chứng chỉ cho Intermediate CA (Hạn dùng 5 năm, pathlen=0)}
\ExtensionTok{openssl}\NormalTok{ x509 }\AttributeTok{{-}req} \AttributeTok{{-}in}\NormalTok{ intermediate.csr }\AttributeTok{{-}CA}\NormalTok{ root\_ca.crt }\AttributeTok{{-}CAkey}\NormalTok{ root\_rsa.key }\DataTypeTok{\textbackslash{}}
  \AttributeTok{{-}CAcreateserial} \AttributeTok{{-}out}\NormalTok{ intermediate.crt }\AttributeTok{{-}days}\NormalTok{ 1825 }\AttributeTok{{-}sha384} \DataTypeTok{\textbackslash{}}
  \AttributeTok{{-}extfile} \OperatorTok{\textless{}(}\BuiltInTok{printf} \StringTok{"basicConstraints=critical,CA:TRUE,pathlen:0\textbackslash{}nkeyUsage=critical,digitalSignature,keyCertSign,cRLSign\textbackslash{}nauthorityKeyIdentifier=keyid:always,issuer\textbackslash{}nsubjectKeyIdentifier=hash"}\OperatorTok{)}

\CommentTok{\# ==============================================================================}
\CommentTok{\# PHẦN 3: CẤP PHÁT CHỨNG CHỈ WEB SERVER VỚI SAN (CA/B FORUM COMPLIANT)}
\CommentTok{\# ==============================================================================}

\CommentTok{\# 9. Tạo CSR cho máy chủ Web với đầy đủ tên miền SAN}
\ExtensionTok{openssl}\NormalTok{ req }\AttributeTok{{-}new} \AttributeTok{{-}newkey}\NormalTok{ rsa:2048 }\AttributeTok{{-}nodes} \AttributeTok{{-}keyout}\NormalTok{ web.key }\AttributeTok{{-}out}\NormalTok{ web.csr }\DataTypeTok{\textbackslash{}}
  \AttributeTok{{-}subj} \StringTok{"/C=VN/ST=Ho Chi Minh/O=Hoc Vien Mat Ma/CN=api.kma.edu.vn"}

\CommentTok{\# 10. Intermediate CA ký cấp chứng chỉ máy chủ Web (Hạn dùng 398 ngày)}
\ExtensionTok{openssl}\NormalTok{ x509 }\AttributeTok{{-}req} \AttributeTok{{-}in}\NormalTok{ web.csr }\AttributeTok{{-}CA}\NormalTok{ intermediate.crt }\AttributeTok{{-}CAkey}\NormalTok{ intermediate.key }\DataTypeTok{\textbackslash{}}
  \AttributeTok{{-}CAcreateserial} \AttributeTok{{-}out}\NormalTok{ web.crt }\AttributeTok{{-}days}\NormalTok{ 398 }\AttributeTok{{-}sha256} \DataTypeTok{\textbackslash{}}
  \AttributeTok{{-}extfile} \OperatorTok{\textless{}(}\BuiltInTok{printf} \StringTok{"basicConstraints=critical,CA:FALSE\textbackslash{}nkeyUsage=digitalSignature,keyEncipherment\textbackslash{}nextendedKeyUsage=serverAuth,clientAuth\textbackslash{}nsubjectAltName=DNS:api.kma.edu.vn,DNS:localhost,IP:127.0.0.1\textbackslash{}nauthorityKeyIdentifier=keyid,issuer"}\OperatorTok{)}

\CommentTok{\# 11. Đóng gói chuỗi tin cậy đầy đủ (Full Chain)}
\FunctionTok{cat}\NormalTok{ web.crt intermediate.crt root\_ca.crt }\OperatorTok{\textgreater{}}\NormalTok{ fullchain.pem}

\CommentTok{\# ==============================================================================}
\CommentTok{\# PHẦN 4: KIỂM ĐỊNH, PHÂN TÍCH NHỊ PHÂN VÀ THU HỒI}
\CommentTok{\# ==============================================================================}

\CommentTok{\# 12. Xem chi tiết văn bản đầy đủ của chứng chỉ X.509}
\ExtensionTok{openssl}\NormalTok{ x509 }\AttributeTok{{-}in}\NormalTok{ web.crt }\AttributeTok{{-}noout} \AttributeTok{{-}text}

\CommentTok{\# 13. Mổ xẻ cấu trúc cây dữ liệu nhị phân ASN.1 DER}
\ExtensionTok{openssl}\NormalTok{ asn1parse }\AttributeTok{{-}in}\NormalTok{ web.crt }\AttributeTok{{-}inform}\NormalTok{ PEM}

\CommentTok{\# 14. Kiểm tra xác thực đường dẫn tin cậy (Chain Validation)}
\ExtensionTok{openssl}\NormalTok{ verify }\AttributeTok{{-}CAfile}\NormalTok{ root\_ca.crt }\AttributeTok{{-}untrusted}\NormalTok{ intermediate.crt web.crt}

\CommentTok{\# 15. Khởi tạo danh sách thu hồi chứng chỉ (CRL) rỗng}
\ExtensionTok{openssl}\NormalTok{ ca }\AttributeTok{{-}gencrl} \AttributeTok{{-}config}\NormalTok{ openssl.cnf }\AttributeTok{{-}out}\NormalTok{ crl.pem}

\CommentTok{\# 16. Chuyển đổi định dạng chứng chỉ từ PEM sang nhị phân DER}
\ExtensionTok{openssl}\NormalTok{ x509 }\AttributeTok{{-}in}\NormalTok{ web.crt }\AttributeTok{{-}outform}\NormalTok{ DER }\AttributeTok{{-}out}\NormalTok{ web.der}

\CommentTok{\# 17. Đóng gói cặp khóa và chứng chỉ sang kho chứa mật khẩu PKCS\#12 (.p12 / .pfx)}
\ExtensionTok{openssl}\NormalTok{ pkcs12 }\AttributeTok{{-}export} \AttributeTok{{-}out}\NormalTok{ bundle.pfx }\AttributeTok{{-}inkey}\NormalTok{ web.key }\AttributeTok{{-}in}\NormalTok{ web.crt }\AttributeTok{{-}certfile}\NormalTok{ intermediate.crt}
\end{Highlighting}
\end{Shaded}

\begin{center}\rule{0.5\linewidth}{0.5pt}\end{center}

\section{PHỤ LỤC B: BỘ CÔNG CỤ PYTHON CRYPTOGRAPHY PHÂN TÍCH \& KIỂM
ĐỊNH}\label{phux1ee5-lux1ee5c-b-bux1ed9-cuxf4ng-cux1ee5-python-cryptography-phuxe2n-tuxedch-kiux1ec3m-ux111ux1ecbnh}

Dưới đây là công cụ Python hoàn chỉnh dựa trên thư viện chuẩn quốc tế
\texttt{cryptography} hỗ trợ học viên tự động hóa kiểm định tính hợp lệ
của chứng chỉ X.509 theo các quy tắc nghiêm ngặt của CA/B Forum:

\begin{Shaded}
\begin{Highlighting}[]
\CommentTok{\#!/usr/bin/env python3}
\CommentTok{\# {-}*{-} coding: utf{-}8 {-}*{-}}
\CommentTok{\textquotesingle{}\textquotesingle{}\textquotesingle{}}
\CommentTok{PKI Certificate Auditor \& RFC 5280 Parser}
\CommentTok{Học phần: Cơ sở hạ tầng khóa công khai (ATHK 512)}
\CommentTok{Tác giả: Nguyễn An Khương (Khoa KH\&KT Máy tính, Trường ĐHBK Tp. HCM)}
\CommentTok{\textquotesingle{}\textquotesingle{}\textquotesingle{}}

\ImportTok{import}\NormalTok{ sys}
\ImportTok{import}\NormalTok{ datetime}
\ImportTok{from}\NormalTok{ cryptography }\ImportTok{import}\NormalTok{ x509}
\ImportTok{from}\NormalTok{ cryptography.hazmat.backends }\ImportTok{import}\NormalTok{ default\_backend}
\ImportTok{from}\NormalTok{ cryptography.x509.oid }\ImportTok{import}\NormalTok{ ExtensionOID, ExtendedKeyUsageOID}

\KeywordTok{def}\NormalTok{ audit\_certificate(cert\_path: }\BuiltInTok{str}\NormalTok{):}
    \BuiltInTok{print}\NormalTok{(}\SpecialStringTok{f"[*] Tiến hành kiểm toán chứng chỉ: }\SpecialCharTok{\{}\NormalTok{cert\_path}\SpecialCharTok{\}}\SpecialStringTok{"}\NormalTok{)}
    \ControlFlowTok{with} \BuiltInTok{open}\NormalTok{(cert\_path, }\StringTok{"rb"}\NormalTok{) }\ImportTok{as}\NormalTok{ f:}
\NormalTok{        data }\OperatorTok{=}\NormalTok{ f.read()}

    \CommentTok{\# Nạp chứng chỉ (hỗ trợ PEM hoặc DER)}
    \ControlFlowTok{try}\NormalTok{:}
\NormalTok{        cert }\OperatorTok{=}\NormalTok{ x509.load\_pem\_x509\_certificate(data, default\_backend())}
    \ControlFlowTok{except} \PreprocessorTok{Exception}\NormalTok{:}
\NormalTok{        cert }\OperatorTok{=}\NormalTok{ x509.load\_der\_x509\_certificate(data, default\_backend())}

\NormalTok{    now }\OperatorTok{=}\NormalTok{ datetime.datetime.now(datetime.timezone.utc)}

    \CommentTok{\# 1. Kiểm tra Version}
\NormalTok{    version\_name }\OperatorTok{=}\NormalTok{ cert.version.name}
    \BuiltInTok{print}\NormalTok{(}\SpecialStringTok{f" [+] Phiên bản X.509: }\SpecialCharTok{\{}\NormalTok{version\_name}\SpecialCharTok{\}}\SpecialStringTok{"}\NormalTok{)}
    \ControlFlowTok{if}\NormalTok{ cert.version }\OperatorTok{!=}\NormalTok{ x509.Version.v3:}
        \BuiltInTok{print}\NormalTok{(}\StringTok{"  [CẢNH BÁO] Chứng chỉ không phải X.509 v3!"}\NormalTok{)}

    \CommentTok{\# 2. Kiểm tra Thời gian hiệu lực}
    \BuiltInTok{print}\NormalTok{(}\SpecialStringTok{f" [+] Hiệu lực từ: }\SpecialCharTok{\{}\NormalTok{cert}\SpecialCharTok{.}\NormalTok{not\_valid\_before\_utc}\SpecialCharTok{\}}\SpecialStringTok{ đến }\SpecialCharTok{\{}\NormalTok{cert}\SpecialCharTok{.}\NormalTok{not\_valid\_after\_utc}\SpecialCharTok{\}}\SpecialStringTok{"}\NormalTok{)}
    \ControlFlowTok{if}\NormalTok{ now }\OperatorTok{\textless{}}\NormalTok{ cert.not\_valid\_before\_utc:}
        \BuiltInTok{print}\NormalTok{(}\StringTok{"  [LỖI] Chứng chỉ chưa đến thời điểm có hiệu lực!"}\NormalTok{)}
    \ControlFlowTok{elif}\NormalTok{ now }\OperatorTok{\textgreater{}}\NormalTok{ cert.not\_valid\_after\_utc:}
        \BuiltInTok{print}\NormalTok{(}\StringTok{"  [LỖI] Chứng chỉ ĐÃ HẾT HẠN!"}\NormalTok{)}
    \ControlFlowTok{else}\NormalTok{:}
        \BuiltInTok{print}\NormalTok{(}\StringTok{"  [PASS] Thời gian hiệu lực hợp lệ."}\NormalTok{)}

    \CommentTok{\# 3. Kiểm tra Độ dài Serial Number (RFC 5280 giới hạn \textless{}= 20 bytes)}
\NormalTok{    serial\_len }\OperatorTok{=}\NormalTok{ (cert.serial\_number.bit\_length() }\OperatorTok{+} \DecValTok{7}\NormalTok{) }\OperatorTok{//} \DecValTok{8}
    \BuiltInTok{print}\NormalTok{(}\SpecialStringTok{f" [+] Serial Number: }\SpecialCharTok{\{}\BuiltInTok{hex}\NormalTok{(cert.serial\_number)}\SpecialCharTok{\}}\SpecialStringTok{ (Độ dài: }\SpecialCharTok{\{}\NormalTok{serial\_len}\SpecialCharTok{\}}\SpecialStringTok{ bytes)"}\NormalTok{)}
    \ControlFlowTok{if}\NormalTok{ serial\_len }\OperatorTok{\textgreater{}} \DecValTok{20}\NormalTok{:}
        \BuiltInTok{print}\NormalTok{(}\StringTok{"  [VI PHẠM RFC 5280] Serial Number vượt quá 20 bytes!"}\NormalTok{)}

    \CommentTok{\# 4. Kiểm tra Basic Constraints}
    \ControlFlowTok{try}\NormalTok{:}
\NormalTok{        bc }\OperatorTok{=}\NormalTok{ cert.extensions.get\_extension\_for\_oid(ExtensionOID.BASIC\_CONSTRAINTS)}
        \BuiltInTok{print}\NormalTok{(}\SpecialStringTok{f" [+] Basic Constraints: is\_ca=}\SpecialCharTok{\{}\NormalTok{bc}\SpecialCharTok{.}\NormalTok{value}\SpecialCharTok{.}\NormalTok{ca}\SpecialCharTok{\}}\SpecialStringTok{, path\_length=}\SpecialCharTok{\{}\NormalTok{bc}\SpecialCharTok{.}\NormalTok{value}\SpecialCharTok{.}\NormalTok{path\_length}\SpecialCharTok{\}}\SpecialStringTok{, critical=}\SpecialCharTok{\{}\NormalTok{bc}\SpecialCharTok{.}\NormalTok{critical}\SpecialCharTok{\}}\SpecialStringTok{"}\NormalTok{)}
        \ControlFlowTok{if}\NormalTok{ bc.value.ca }\KeywordTok{and} \KeywordTok{not}\NormalTok{ bc.critical:}
            \BuiltInTok{print}\NormalTok{(}\StringTok{"  [CẢNH BÁO] Chứng chỉ CA nhưng cờ Basic Constraints không đặt CRITICAL!"}\NormalTok{)}
    \ControlFlowTok{except}\NormalTok{ x509.ExtensionNotFound:}
        \BuiltInTok{print}\NormalTok{(}\StringTok{"  [LỖI NGHIÊM TRỌNG] Thiếu phần mở rộng Basic Constraints!"}\NormalTok{)}

    \CommentTok{\# 5. Kiểm tra Subject Alternative Name (SAN)}
    \ControlFlowTok{try}\NormalTok{:}
\NormalTok{        san }\OperatorTok{=}\NormalTok{ cert.extensions.get\_extension\_for\_oid(ExtensionOID.SUBJECT\_ALTERNATIVE\_NAME)}
\NormalTok{        names }\OperatorTok{=}\NormalTok{ san.value.get\_values\_for(x509.DNSName)}
\NormalTok{        ips }\OperatorTok{=}\NormalTok{ [}\BuiltInTok{str}\NormalTok{(ip) }\ControlFlowTok{for}\NormalTok{ ip }\KeywordTok{in}\NormalTok{ san.value.get\_values\_for(x509.IPAddress)]}
        \BuiltInTok{print}\NormalTok{(}\SpecialStringTok{f" [+] SAN DNS Names: }\SpecialCharTok{\{}\NormalTok{names}\SpecialCharTok{\}}\SpecialStringTok{"}\NormalTok{)}
        \BuiltInTok{print}\NormalTok{(}\SpecialStringTok{f" [+] SAN IP Addresses: }\SpecialCharTok{\{}\NormalTok{ips}\SpecialCharTok{\}}\SpecialStringTok{"}\NormalTok{)}
    \ControlFlowTok{except}\NormalTok{ x509.ExtensionNotFound:}
        \BuiltInTok{print}\NormalTok{(}\StringTok{"  [VI PHẠM CA/B FORUM] Thiếu phần mở rộng Subject Alternative Name (SAN)!"}\NormalTok{)}

    \BuiltInTok{print}\NormalTok{(}\StringTok{"[*] Kiểm toán hoàn tất."}\NormalTok{)}

\ControlFlowTok{if} \VariableTok{\_\_name\_\_} \OperatorTok{==} \StringTok{"\_\_main\_\_"}\NormalTok{:}
    \ControlFlowTok{if} \BuiltInTok{len}\NormalTok{(sys.argv) }\OperatorTok{\textless{}} \DecValTok{2}\NormalTok{:}
        \BuiltInTok{print}\NormalTok{(}\SpecialStringTok{f"Cách dùng: python3 }\SpecialCharTok{\{}\NormalTok{sys}\SpecialCharTok{.}\NormalTok{argv[}\DecValTok{0}\NormalTok{]}\SpecialCharTok{\}}\SpecialStringTok{ \textless{}duong\_dan\_file\_chung\_chi\textgreater{}"}\NormalTok{)}
\NormalTok{        sys.exit(}\DecValTok{1}\NormalTok{)}
\NormalTok{    audit\_certificate(sys.argv[}\DecValTok{1}\NormalTok{])}
\end{Highlighting}
\end{Shaded}

\begin{center}\rule{0.5\linewidth}{0.5pt}\end{center}

\section{PHỤ LỤC C: DANH MỤC TÀI LIỆU THAM KHẢO \& KHÓA HỌC QUỐC
TẾ}\label{phux1ee5-lux1ee5c-c-danh-mux1ee5c-tuxe0i-liux1ec7u-tham-khux1ea3o-khuxf3a-hux1ecdc-quux1ed1c-tux1ebf}

\subsubsection{1. Hệ Thống Khóa Học Quốc Tế Trực Tuyến Hàng Đầu
(OpenCourseWare \& Video
Lectures)}\label{hux1ec7-thux1ed1ng-khuxf3a-hux1ecdc-quux1ed1c-tux1ebf-trux1ef1c-tuyux1ebfn-huxe0ng-ux111ux1ea7u-opencourseware-video-lectures}

\begin{itemize}
\tightlist
\item
  \textbf{MIT OpenCourseWare 6.857: Computer and Network Security}\\
  \emph{Giảng viên:} GS. Ronald L. Rivest (Đồng tác giả RSA) và GS.
  Shafi Goldwasser (Giải thưởng Turing).\\
  \emph{Tài nguyên:} Bài giảng lý thuyết số, định danh chứng thực số và
  lược đồ cam kết mật mã.\\
  \emph{Địa chỉ truy cập:}
  \href{https://ocw.mit.edu/courses/6-857-computer-and-network-security-spring-2014/}{ocw.mit.edu/courses/6-857-computer-and-network-security-spring-2014}
\item
  \textbf{Stanford University CS 255: Introduction to Cryptography}\\
  \emph{Giảng viên:} GS. Dan Boneh.\\
  \emph{Nội dung:} Phân tích chuyên sâu về hệ mật khóa công khai, tính
  toán cặp trên đường cong elliptic (Pairing-based Crypto), thuật toán
  chữ ký số RSA-PSS và ECDSA.\\
  \emph{Địa chỉ truy cập:}
  \href{https://crypto.stanford.edu/~dabo/cs255/}{crypto.stanford.edu/\textasciitilde dabo/cs255/}
\item
  \textbf{UC Berkeley CS 161: Computer Security}\\
  \emph{Giảng viên:} GS. Nicholas Weaver và GS. David Wagner.\\
  \emph{Nội dung:} Thực hành tấn công và phòng thủ giao thức TLS/SSL,
  phân tích lỗ hổng CA, mô hình mối đe dọa Man-in-the-Middle và chính
  sách bảo mật trình duyệt.\\
  \emph{Địa chỉ truy cập:} \href{https://cs161.org/}{cs161.org}
\item
  \textbf{Simons Institute for the Theory of Computing (UC Berkeley):}\\
  \emph{Chuyên đề:} \emph{Lattices: Algorithms, Complexity, and
  Cryptography} \& \emph{Quantum Cryptography Program}.\\
  \emph{Địa chỉ:}
  \href{https://simons.berkeley.edu/}{simons.berkeley.edu}
\item
  \textbf{IACR Real World Crypto (RWC) Symposium:}\\
  \emph{Kênh tài liệu video:} Hội nghị thường niên quốc tế quy tụ các kỹ
  sư hàng đầu từ Google, Cloudflare, Apple về thực nghiệm triển khai TLS
  1.3, PQC Migration và tấn công kênh bên vào HSM.\\
  \emph{Kênh lưu trữ:} \href{https://rwc.iacr.org/}{rwc.iacr.org}
\end{itemize}

\subsubsection{2. Tiêu Chuẩn Kỹ Thuật Quốc Tế Bắt Buộc (Standards \&
RFCs)}\label{tiuxeau-chuux1ea9n-kux1ef9-thuux1eadt-quux1ed1c-tux1ebf-bux1eaft-buux1ed9c-standards-rfcs}

\begin{enumerate}
\def\labelenumi{\arabic{enumi}.}
\tightlist
\item
  \textbf{IETF RFC 5280:} \emph{Internet X.509 Public Key Infrastructure
  Certificate and Certificate Revocation List (CRL) Profile}, 2008.
\item
  \textbf{IETF RFC 3647:} \emph{Internet X.509 Public Key Infrastructure
  Certificate Policy and Certification Practices Framework}, 2003.
\item
  \textbf{IETF RFC 8555:} \emph{Automatic Certificate Management
  Environment (ACME)}, 2019.
\item
  \textbf{IETF RFC 6960:} \emph{X.509 Internet Public Key Infrastructure
  Online Certificate Status Protocol - OCSP}, 2013.
\item
  \textbf{IETF RFC 6962:} \emph{Certificate Transparency}, 2013.
\item
  \textbf{IETF RFC 8446:} \emph{The Transport Layer Security (TLS)
  Protocol Version 1.3}, 2018.
\item
  \textbf{NIST FIPS 203:} \emph{Module-Lattice-Based Key-Encapsulation
  Mechanism Standard (ML-KEM)}, 2024.
\item
  \textbf{NIST FIPS 204:} \emph{Module-Lattice-Based Digital Signature
  Standard (ML-DSA)}, 2024.
\item
  \textbf{NIST FIPS 205:} \emph{Stateless Hash-Based Digital Signature
  Standard (SLH-DSA)}, 2024.
\item
  \textbf{CA/Browser Forum:} \emph{Baseline Requirements for the
  Issuance and Management of Publicly-Trusted Certificates}, Version
  2.0+, 2024.
\end{enumerate}

\end{document}
